\begin{derivation}[从\eqnrefs{eq:sm-top-bw-amplitude-dp94}到\eqnrefs{eq:sm-top-helicity-fractions-dp94}]
在顶夸克静止系取
\begin{equation*}
 p_t^\mu=(m_tc,0,0,0),
 \qquad
 k^\mu=\left(E_W/c,0,0,+|\mathbf k|\right),
 \qquad
 p_b^\mu=\left(E_b/c,0,0,-|\mathbf k|\right),
\end{equation*}
且 $m_b=0$。二体运动学给
\begin{equation*}
 E_W=\frac{m_t^2+m_W^2}{2m_t}c^2,
 \qquad
 E_b=|\mathbf k|c=\frac{m_t^2-m_W^2}{2m_t}c^2.
\end{equation*}
对 $W$ 取极化矢量
\begin{align*}
 \varepsilon_0^\mu
 &=\frac1{m_Wc}\left(|\mathbf k|,0,0,E_W/c\right),\\
 \varepsilon_\pm^\mu
 &=\frac1{\sqrt2}(0,\mp1,-\ii,0),
\end{align*}
满足 $k\cdot\varepsilon_\lambda=0$ 与 $\varepsilon_\lambda^*\cdot\varepsilon_\lambda=-1$。

固定 $W$ 极化后，对末态 $b$ 自旋求和并对初态顶夸克自旋平均。利用
\(\sum_su_s(p)\bar u_s(p)=\slashed p+mc\)，有
\begin{equation*}
 \overline{|\mathcal M_\lambda|^2}
 =\frac{g^2|V_{tb}|^2}{4}
 \varepsilon_{\mu}^{(\lambda)*}\varepsilon_\nu^{(\lambda)}
 L^{\mu\nu},
\qquad
 L^{\mu\nu}\equiv
 \Tr\!\left[
 \slashed p_b\gamma^\mu P_L(\slashed p_t+m_tc)
 \gamma^\nu P_L
 \right].
\end{equation*}
先把这个迹本身算清楚。由于
\(P_L\gamma^\nu=\gamma^\nu P_R\) 且 $P_RP_L=0$，质量项满足
\begin{equation*}
 m_tc\,\Tr(\slashed p_b\gamma^\mu P_L\gamma^\nu P_L)=0.
\end{equation*}
所以只剩四动量项。展开两个手征投影，含单个 $\gamma^5$ 的部分用
\(\Tr(\gamma^\alpha\gamma^\mu\gamma^\beta\gamma^\nu\gamma^5)
=-4\ii\epsilon^{\alpha\mu\beta\nu}\)，不含 $\gamma^5$ 的部分用四伽马迹，得到
\begin{equation*}
\boxed{
L^{\mu\nu}
=2\left[
 p_b^\mu p_t^\nu+p_b^\nu p_t^\mu
 -\eta^{\mu\nu}(p_b\cdot p_t)
 -\ii\epsilon^{\mu\nu\alpha\beta}p_{b\alpha}p_{t\beta}
 \right].}
\end{equation*}
这一步把“$V-A$”的两种数学作用分开了：前三项来自矢量迹，最后的反对称项来自 $\gamma^5$，它正负责区分两个横向螺旋度。

在顶夸克静止系中记
\begin{equation*}
 Q\equiv p_t\cdot p_b
 =\frac{m_t^2c^2}{2}(1-x_W).
\end{equation*}
代入上面的 $p_t,p_b$ 后，$L^{\mu\nu}$ 的非零部分可直接写成
\begin{equation*}
L^{\mu\nu}
=2Q
\begin{pmatrix}
 1&0&0&-1\\
 0&1&\ii&0\\
 0&-\ii&1&0\\
 -1&0&0&1
\end{pmatrix}^{\mu\nu}.
\end{equation*}
现在三个螺旋度无需再借助口头角动量论证，而可以逐个收缩。

对正螺旋度，协变偏振分量为
\begin{equation*}
\varepsilon_{\mu}^{(+)}
=\frac1{\sqrt2}(0,1,\ii,0).
\end{equation*}
横向 $x$--$y$ 子块作用在列向量 $(1,\ii)^{\mathsf T}$ 上恰好为零：
\begin{equation*}
\begin{pmatrix}1&\ii\\-\ii&1\end{pmatrix}
\begin{pmatrix}1\\\ii\end{pmatrix}
=\begin{pmatrix}0\\0\end{pmatrix}.
\end{equation*}
因此
\begin{equation*}
\boxed{\overline{|\mathcal M_+|^2}=0}.
\end{equation*}
这里的零来自手征张量中对称部分与 $\gamma^5$ 反对称部分的精确相消；这就是 $m_b=0$ 时右手 $W$ 被禁止的代数形式。

对负螺旋度，
\(\varepsilon_{\mu}^{(-)}=(0,-1,\ii,0)/\sqrt2\)，于是
\begin{align*}
\varepsilon_{\mu}^{(-)*}L^{\mu\nu}\varepsilon_{\nu}^{(-)}
&=4Q
=2m_t^2c^2(1-x_W).
\end{align*}
按正文采用的约化振幅归一化，这给出
\begin{equation*}
 \boxed{
 \overline{|\mathcal M_-|^2}
 =\frac{g^2|V_{tb}|^2}{2}
 m_t^2(1-x_W)}.
\end{equation*}

纵向偏振只有 $0,z$ 分量。令
\begin{equation*}
 a\equiv\frac{|\mathbf k|}{m_Wc},
 \qquad
 b\equiv\frac{E_W}{m_Wc^2},
 \qquad
 \varepsilon_{\mu}^{(0)}=(a,0,0,-b).
\end{equation*}
$0$--$z$ 子块同样可以直接收缩：
\begin{align*}
\varepsilon_{\mu}^{(0)}L^{\mu\nu}\varepsilon_{\nu}^{(0)}
&=2Q(a+b)^2.
\end{align*}
而二体运动学给
\begin{align*}
 a+b
&=\frac{|\mathbf k|c+E_W}{m_Wc^2}\\
&=\frac{m_tc^2(1-x_W)/2+m_tc^2(1+x_W)/2}{m_Wc^2}\\
&=\frac{m_t}{m_W}.
\end{align*}
因此
\begin{equation*}
\varepsilon_{\mu}^{(0)*}L^{\mu\nu}\varepsilon_{\nu}^{(0)}
=\frac{2Q}{x_W},
\end{equation*}
从而
\begin{equation*}
 \boxed{
 \overline{|\mathcal M_0|^2}
 =\frac{g^2|V_{tb}|^2}{4}
 \frac{m_t^4}{m_W^2}(1-x_W)}.
\end{equation*}
纵向增强的来源现在是显式的：$\varepsilon_0^\mu$ 自身含 $1/m_W$，而 $a+b=m_t/m_W$，所以当 $m_t\gg m_W$ 时纵向通道相对于横向通道增强。
因此
\begin{equation*}
 \frac{\overline{|\mathcal M_-|^2}}
 {\overline{|\mathcal M_0|^2}}
 =2\frac{m_W^2}{m_t^2}=2x_W,
\end{equation*}
而 $\overline{|\mathcal M_+|^2}=0$。用
\begin{equation*}
F_\lambda
=\frac{\overline{|\mathcal M_\lambda|^2}}
{\overline{|\mathcal M_0|^2}+\overline{|\mathcal M_-|^2}+\overline{|\mathcal M_+|^2}}
\end{equation*}
逐项除以 $\overline{|\mathcal M_0|^2}$，得到
\begin{equation*}
 F_0=\frac1{1+2x_W},
 \qquad
 F_-=\frac{2x_W}{1+2x_W},
 \qquad
 F_+=0.
\end{equation*}

总自旋平均振幅平方为
\begin{equation*}
 \overline{|\mathcal M|^2}
 =\frac{g^2|V_{tb}|^2}{4}
 \frac{m_t^4}{m_W^2}(1-x_W)(1+2x_W).
\end{equation*}
二体衰变公式
\begin{equation*}
 \Gamma
 =\frac{|\mathbf k|}{8\pi\hbar m_t^2c}
 \overline{|\mathcal M|^2}
\end{equation*}
与 $|\mathbf k|=m_tc(1-x_W)/2$ 合并，再用
\begin{equation*}
 \frac{G_F^{(E)}}{\sqrt2}=\frac{g^2}{8(m_Wc^2)^2},
\end{equation*}
即得到正文\eqnrefs{eq:sm-top-width-dp94}。
\end{derivation}

\begin{derivation}[从\eqnrefs{eq:sm-gauge-lagrangian-dp11}到\eqnrefs{eq:sm-wwv-vertex-dp11}]
一般 Yang--Mills 三点张量由场强的非线性项固定。对 $SU(2)_L$，生成元满足

\[
 [t^a,t^b]=\ii\epsilon^{abc}t^c,
\]
所以结构常数就是 $f^{abc}=\epsilon^{abc}$。由于 $\epsilon^{abc}$ 在任意两个指标相同时为零，含一对 $W^+,W^-$ 的非零三点项必须由原始 $\mathcal W^1,\mathcal W^2$ 分量组成，而第三条原始中性规范场是 $\mathcal W^3$。把带电场定义的逆变换
\begin{equation*}
 \mathcal W^1=\frac{\mathcal W^++\mathcal W^-}{\sqrt2},
 \qquad
 \mathcal W^2=\frac{\ii(\mathcal W^+-\mathcal W^-)}{\sqrt2}
\end{equation*}
代入群因子。固定第一条外腿为 $W^+$、第二条为 $W^-$、第三条为 $\mathcal W^3$ 时，两个非零实基分量给
\begin{align*}
 \sum_{a,b=1}^{2}C_a^+C_b^-\epsilon^{ab3}
 &=\frac1{\sqrt2}\left(-\frac{\ii}{\sqrt2}\right)\epsilon^{123}
 +\frac{\ii}{\sqrt2}\frac1{\sqrt2}\epsilon^{213}\\
 &=-\frac{\ii}{2}-\frac{\ii}{2}\\
 &=-\ii.
\end{align*}
一般实基顶角的群因子是 $-g\epsilon^{ab3}$，所以带电基中 $W^+W^-\mathcal W^3$ 的总体因子变为 $+\ii g$。洛伦兹张量仍保持
\begin{equation*}
 \mathcal V_{\mu\nu\rho}(p,q,r)
 =\eta_{\mu\nu}(p-q)_\rho
 +\eta_{\nu\rho}(q-r)_\mu
 +\eta_{\rho\mu}(r-p)_\nu.
\end{equation*}
再用
\begin{equation*}
 \mathcal W^3=s_WA+c_W\mathcal Z
\end{equation*}
得到
\begin{equation*}
 g\mathcal W^3=gs_WA+gc_W\mathcal Z.
\end{equation*}
由\eqnrefs{eq:gem-definition-dp11} 的 $gs_W=g_{\rm em}$ 代入上式，得到
\begin{equation*}
 g_{WW\gamma}=g_{\rm em},
 \qquad
 g_{WWZ}=gc_W,
\end{equation*}
即\eqnrefs{eq:sm-wwv-couplings-dp11}。把同一 $\mathcal V_{\mu\nu\rho}$ 的三条动量依次指定为 $p_+,p_-,k$，并保持全部流入约定，便得到\eqnrefs{eq:sm-wwv-vertex-dp11}；其中显式的 $+\ii$ 正来自上述实基到带电基的复线性变换。
\end{derivation}

\begin{derivation}[从\eqnrefs{eq:sm-wwv-vertex-dp11}到\eqnrefs{eq:sm-wwgamma-ward-dp41}]
正文\eqnrefs{eq:sm-wwv-vertex-dp11} 已固定三矢量顶角的动量结构。取 $V=\gamma$ 并用光子四动量收缩，再借助三条外腿全部流入的动量守恒，就能把结果整理成两条 $W$ 逆传播子之差，即正文\eqnrefs{eq:sm-wwgamma-ward-dp41}。

从\eqnrefs{eq:sm-wwv-vertex-dp11} 取 $V=\gamma$ 并与光子动量 $k^\rho$ 收缩：
\begin{align*}
 \frac{k^\rho\Gamma^{WW\gamma}_{\mu\nu\rho}}{+\ii g_{\rm em}}
 ={}&\eta_{\mu\nu}k\cdot(p_+-p_-)
 +k_\nu(p_--k)_\mu
 +k_\mu(k-p_+)_\nu.
\end{align*}
动量守恒 $k=-p_+-p_-$ 给出第一项
\begin{align*}
 k\cdot(p_+-p_-)
 &=(-p_+-p_-)\cdot(p_+-p_-)\\
 &=-p_+^2+p_-^2.
\end{align*}
后两项分别展开为
\begin{align*}
 k_\nu(p_--k)_\mu
 &=-(p_{+\nu}+p_{-\nu})(p_{+\mu}+2p_{-\mu}),\\
 k_\mu(k-p_+)_\nu
 &=(p_{+\mu}+p_{-\mu})(2p_{+\nu}+p_{-\nu}).
\end{align*}
把乘积逐项展开，$p_{+\mu}p_{-\nu}$ 与 $p_{-\mu}p_{+\nu}$ 的交叉项互相抵消，剩下
\begin{equation*}
 k_\nu(p_--k)_\mu+k_\mu(k-p_+)_\nu
 =p_{+\mu}p_{+\nu}-p_{-\mu}p_{-\nu}.
\end{equation*}
因此
\begin{align*}
 k^\rho\Gamma^{WW\gamma}_{\mu\nu\rho}
 =+\ii g_{\rm em}\Big[
 &(p_-^2\eta_{\mu\nu}-p_{-\mu}p_{-\nu})\\
 &-(p_+^2\eta_{\mu\nu}-p_{+\mu}p_{+\nu})
 \Big].
\end{align*}
在两个括号中同时减去相同的质量项 $m_W^2c^2\eta_{\mu\nu}$ 不改变差值，于是恰好得到
\begin{equation*}
 k^\rho\Gamma^{WW\gamma}_{\mu\nu\rho}
 =+\ii g_{\rm em}
 \left[\Pi_{\mu\nu}(p_-)-\Pi_{\mu\nu}(p_+)\right],
\end{equation*}
即\eqnrefs{eq:sm-wwgamma-ward-dp41}。若两条 $W$ 外腿在壳，则
\begin{equation*}
 p_\pm^2=m_W^2c^2,
 \qquad
 p_\pm\cdot\varepsilon_\pm=0,
\end{equation*}
所以 $\varepsilon_+^\mu\Pi_{\mu\nu}(p_+)=0$ 且 $\Pi_{\mu\nu}(p_-)\varepsilon_-^\nu=0$。这给出外部光子纵向替换后的零振幅，并同时检查了动量号、质量项与偏振条件。
\end{derivation}

\begin{derivation}[质量因子、自旋迹与二体相空间]
使用 $P_E=L_E+N_E$、$\bar u(L_E)\slashed L_E=M_\ell\bar u(L_E)$ 和 $\slashed N_Ev(N_E)=0$，
\begin{align*}
 P_E^\mu\bar u\gamma_\mu(1-\gamma^5)v
 &=\bar u(\slashed L_E+\slashed N_E)(1-\gamma^5)v\\
 &=M_\ell\bar u(1-\gamma^5)v
 +\bar u(1+\gamma^5)\slashed N_Ev\\
 &=M_\ell\bar u(1-\gamma^5)v.
\end{align*}
因此抑制不是通过先比较两个相空间体积猜出的。对末态自旋求和，初态为自旋零而无自旋平均，
\begin{align*}
 \sum_{\rm spins}|\bar u(1-\gamma^5)v|^2
 &=\Tr[(\slashed L_E+M_\ell)(1-\gamma^5)
 \slashed N_E(1+\gamma^5)]\\
 &=2\Tr[(\slashed L_E+M_\ell)\slashed N_E(1+\gamma^5)]\\
 &=8L_E\cdot N_E=4(M_P^2-M_\ell^2),\\
 \sum|\mathcal M_{P\ell2}|^2
 &=2[G_F^{(E)}]^2|V_{ij}|^2f_P^2M_\ell^2(M_P^2-M_\ell^2).
\end{align*}
能量四矢量归一化下，二体相空间积分为
\begin{equation*}
 \int\dd\Phi_{2,E}=\frac{1}{8\pi}\left(1-\frac{M_\ell^2}{M_P^2}\right).
\end{equation*}
再用 $\Gamma=(2\hbar M_P)^{-1}\int\dd\Phi_{2,E}\sum|\mathcal M|^2$，即得\eqnrefs{eq:sm-pseudoscalar-width}。一个阈值因子来自自旋迹，另一个来自相空间，这解释了最终的平方。
\end{derivation}

\begin{derivation}[从\eqnrefs{eq:hggamma-transverse-tensor-dp11}到\eqnrefs{eq:sm-hgamma-loop-functions-dp41}]
由 Yukawa 相互作用，$h f\bar f$ 的约化顶角为 $-\ii m_fc^2/v_E$，光子顶角为 $-\ii g_{\rm em}Q_f\gamma^\mu$。 两种相反的光子排序对应两张三角图，必须同时保留。 提出共同的外腿归一化因子后，定义约化张量

\begin{align*}
 \mathcal M_f^{\mu\nu}
 ={}&-N_c^f(g_{\rm em}Q_f)^2\frac{m_fc^2}{v_E}
 \int_\ell^{(d)}
 \frac{\Tr\!\left[
 (\slashed\ell+m_fc)\gamma^\mu
 (\slashed\ell+\slashed k_1+m_fc)\gamma^\nu
 (\slashed\ell-\slashed k_2+m_fc)
 \right]}
 {D_0D_1D_2}
 \nonumber\\
 &+(\mu,k_1\leftrightarrow\nu,k_2),
\end{align*}
其中
\begin{align*}
 D_0&=\ell^2-m_f^2c^2+\ii0,\\
 D_1&=(\ell+k_1)^2-m_f^2c^2+\ii0,\\
 D_2&=(\ell-k_2)^2-m_f^2c^2+\ii0,
\end{align*}
两个外部光子满足
\begin{equation*}
 k_1^2=k_2^2=0,
 \qquad
 (k_1+k_2)^2=m_h^2c^2,
 \qquad
 2k_1\!\cdot k_2=m_h^2c^2.
\end{equation*}
闭合费米子环给出该振幅的整体负号。

三个分母用单纯形上的 Feynman 参数合并：
\begin{equation*}
 \frac1{D_0D_1D_2}
 =2\int_0^1\dd x\int_0^{1-x}\dd y\,
 \frac1{[xD_1+yD_2+(1-x-y)D_0]^3}.
\end{equation*}
逐项展开括号：
\begin{align*}
 xD_1+yD_2+(1-x-y)D_0
 ={}&\ell^2+2\ell\!\cdot(xk_1-yk_2)
 -m_f^2c^2
 \\
 &-2xy\,k_1\!\cdot k_2+\ii0.
\end{align*}
令
\begin{equation*}
 L\equiv\ell+xk_1-yk_2,
\end{equation*}
并用 $k_1^2=k_2^2=0$，分母化成
\begin{equation*}
 [L^2-\Delta_f(x,y)+\ii0]^3,
 \qquad
 {\Delta_f=m_f^2c^2-xy\,m_h^2c^2.}
\end{equation*}

接着把分子真正展开。为避免在长迹中反复书写，记
\[
a\equiv\ell,\qquad b\equiv\ell+k_1,\qquad c\equiv\ell-k_2,
\qquad M_f\equiv m_fc.
\]
第一种光子排序的迹为
\begin{align*}
N^{\mu\nu}
&\equiv\Tr\!\left[(\slashed a+M_f)\gamma^\mu
(\slashed b+M_f)\gamma^\nu(\slashed c+M_f)\right].
\end{align*}
含零个或两个 $M_f$ 的项分别含奇数个 $\gamma$ 矩阵，迹为零；只剩一个 $M_f$ 和三个 $M_f$ 的项。逐项使用
\[
\Tr(\gamma^\alpha\gamma^\beta)=4\eta^{\alpha\beta},\qquad
\Tr(\gamma^\alpha\gamma^\beta\gamma^\gamma\gamma^\delta)
=4(\eta^{\alpha\beta}\eta^{\gamma\delta}
-\eta^{\alpha\gamma}\eta^{\beta\delta}
+\eta^{\alpha\delta}\eta^{\beta\gamma}),
\]
得到
\begin{align*}
\frac{N^{\mu\nu}}{4M_f}
={}&b^\mu c^\nu+c^\mu b^\nu-\eta^{\mu\nu}b\!\cdot c
+a^\mu c^\nu-a^\nu c^\mu+\eta^{\mu\nu}a\!\cdot c\\
&+a^\mu b^\nu+a^\nu b^\mu-\eta^{\mu\nu}a\!\cdot b
+M_f^2\eta^{\mu\nu}.
\end{align*}
第二张三角图由 $(\mu,k_1)\leftrightarrow(\nu,k_2)$ 得到。于是两张图相加以后，整个问题已经从六个 $\gamma$ 矩阵的迹化成普通 Lorentz 张量代数。

平移到 $L=\ell+xk_1-yk_2$ 后，奇数次 $L$ 的积分为零；秩二项用
\begin{equation*}
\int\dd^dL\,L^\mu L^\nu F(L^2)
=\frac{\eta^{\mu\nu}}{d}\int\dd^dL\,L^2F(L^2)
\end{equation*}
约化。为了不在中间步骤预先假定横向性，定义正文张量
\[
T^{\mu\nu}\equiv k_1\!\cdot k_2\,\eta^{\mu\nu}-k_2^\mu k_1^\nu,
\]
其自收缩满足
\begin{equation*}
T_{\mu\nu}T^{\mu\nu}=(d-2)(k_1\!\cdot k_2)^2.
\end{equation*}
因此横向系数可以在所有张量积分完成后由
\begin{equation*}
\mathcal C_f
=\frac{T_{\mu\nu}\mathcal M_f^{\mu\nu}}
{(d-2)(k_1\!\cdot k_2)^2}
\end{equation*}
直接投影。把上面的显式 $N^{\mu\nu}$ 代入这一投影，所有 $L^2$ 项都通过
\begin{equation*}
\frac{L^2}{(L^2-\Delta_f)^3}
=\frac1{(L^2-\Delta_f)^2}
+\frac{\Delta_f}{(L^2-\Delta_f)^3}
\end{equation*}
化成二点型与三点型标量积分；两种光子排序相加后，二点型局域项彼此抵消。剩余三点项为
\begin{equation*}
 A_{1/2}(\tau)
 =4\tau\int_0^1\dd x\int_0^{1-x}\dd y\,
 \frac{1-4xy}{\tau-4xy-\ii0},
 \qquad
 \tau\equiv\frac{4m_f^2}{m_h^2}.
\end{equation*}
这里的分子 $1-4xy$ 是费米子迹中质量项与动量项共同作用的结果；只保留其中任一部分都会破坏重质量极限。

双积分可用下列组合化为闭式：
\begin{equation*}
 I(\tau)
 \equiv\int_0^1\dd x\int_0^{1-x}\dd y\,
 \frac1{\tau-4xy-\ii0}.
\end{equation*}
先对 $y$ 积分：
\begin{align*}
 I(\tau)
 &=\int_0^1\frac{\dd x}{4x}
 \ln\!\left[
 \frac{\tau}{\tau-4x(1-x)-\ii0}
 \right].
\end{align*}
再令 $x=(1-z)/2$，使 $4x(1-x)=1-z^2$。对 $\tau\ge1$，积分路径不穿过分支割线；完成变量代换并积分后得到
\begin{equation*}
 I(\tau)=\frac12\arcsin^2\!\frac1{\sqrt\tau}.
\end{equation*}
对 $\tau<1$，分母在积分区域内穿过双费米子阈值；按 $-\ii0$ 处方解析延拓后出现对数与虚部。因此统一定义
\begin{equation*}
 {
 f(\tau)=2I(\tau)=
 \begin{cases}
 \arcsin^2(1/\sqrt\tau),&\tau\ge1,\\[4pt]
 -\dfrac14\left[
 \ln\dfrac{1+\sqrt{1-\tau}}{1-\sqrt{1-\tau}}-\ii\pi
 \right]^2,&\tau<1.
 \end{cases}}
\end{equation*}
$\tau=1$ 正是 $m_h=2m_f$ 的粒子对产生分支点。

现在把上式的分子写成
\begin{equation*}
 1-4xy=(\tau-4xy)+(1-\tau).
\end{equation*}
于是
\begin{align*}
 A_{1/2}(\tau)
 &=4\tau\left[
 \int_\triangle\dd x\dd y
 +(1-\tau)I(\tau)
 \right],\\
 &=4\tau\left[\frac12+\frac{1-\tau}{2}f(\tau)\right],\\
 &={2\tau[1+(1-\tau)f(\tau)]}.
\end{align*}
因此 $A_{1/2}(\tau)$ 恰好化为\eqnrefs{eq:sm-hgamma-loop-functions-dp41} 的第一行。

重费米子极限也可以直接从参数积分检验。$\tau\gg1$ 时
\begin{equation*}
 f(\tau)=\frac1\tau+\frac1{3\tau^2}+\order(\tau^{-3}),
\end{equation*}
代回得到
\begin{equation*}
 {A_{1/2}(\tau)=\frac43+\order(\tau^{-1}).}
\end{equation*}
\end{derivation}

\begin{derivation}[从\eqnrefs{eq:hggamma-transverse-tensor-dp11}到\eqnrefs{eq:sm-hgamma-loop-functions-dp41}]
$W$ 环的困难不在最后一个参数积分，而在自旋一传播子的纵向分子会把圈动量幂次显著抬高。采用幺正规范后，Goldstone 与 ghost 不再作为独立内部线出现，而物理在壳振幅保持规范无关。记
\[
M_W^2\equiv m_W^2c^2,\qquad M_h^2\equiv m_h^2c^2,
\]
则 $W$ 传播子与 $hW^+W^-$ 约化顶角分别为
\begin{equation*}
 D_W^{\alpha\beta}(q)
 =\frac{-\ii}{q^2-M_W^2+\ii0}
 \left(\eta^{\alpha\beta}-\frac{q^\alpha q^\beta}{M_W^2}\right),
 \qquad
 V_{hWW}^{\alpha\beta}=\ii\frac{2(m_Wc^2)^2}{v_E}\eta^{\alpha\beta}.
\end{equation*}
两个 $WW\gamma$ 顶角由正文已建立的 Yang--Mills 顶角给出，另有一个 $WW\gamma\gamma$ 接触顶角。因此幺正规范下需要三角图、交换两个光子的交叉三角图以及接触图，缺少其中任一项都会破坏光子 Ward 恒等式。

令两个光子均出射，$k_1^2=k_2^2=0$，$2k_1\!\cdot k_2=M_h^2$；取三角图的三个分母
\begin{equation*}
D_0=p^2-M_W^2+\ii0,\qquad
D_1=(p-k_1)^2-M_W^2+\ii0,\qquad
D_2=(p-k_1-k_2)^2-M_W^2+\ii0.
\end{equation*}
提出共同因子
\begin{equation*}
\mathcal N_W\equiv\frac{g_{\rm em}^2g}{16\pi^2m_Wc},
\end{equation*}
并暂时不收缩外部偏振。把三张图的 Lorentz 指标全部收缩后，所有含 $k_1^\mu$ 或 $k_2^\nu$ 的项在物理外腿上消失；其余部分可以在五个张量上展开为
\begin{equation*}
\mathcal M_W^{\mu\nu}
=\int_p^{(d)}\!
\left[
\mathscr M_1\eta^{\mu\nu}
+\mathscr M_2p^\mu p^\nu
+\mathscr M_3p^\mu k_1^\nu
+\mathscr M_4k_2^\mu p^\nu
+\mathscr M_5k_2^\mu k_1^\nu
\right].
\end{equation*}
五个系数先分别完成圈积分与张量约化；把同一规范不变图集合求和后，它们自动重组为满足 Ward 恒等式的横向张量。

按正文\eqnrefs{eq:pv-scalar-masters-dp90,eq:pv-rank2-decomposition-dp90} 做 Passarino--Veltman 约化。为缩短公式，以下记
\begin{equation*}
 A_0\equiv A_0(M_W^2),\qquad
 B_h\equiv B_0(M_h^2;M_W^2,M_W^2),\qquad
 B_0\equiv B_0(0;M_W^2,M_W^2),
\end{equation*}
\begin{equation*}
 C_h\equiv C_0(M_h^2,0,0;M_W^2,M_W^2,M_W^2).
\end{equation*}
五类张量积分分别化为
\begin{align*}
\int_p^{(d)}\mathscr M_1\eta^{\mu\nu}
={}&\mathcal N_W
\Bigl[4M_W^2(1-2M_h^2C_h)-(M_h^2+6M_W^2)B_h\Bigr]\eta^{\mu\nu},\\[2mm]
\int_p^{(d)}\mathscr M_2p^\mu p^\nu
={}&\frac{\mathcal N_W}{M_h^2}
\Bigl\{
(M_h^2+6M_W^2)(1+B_h+2M_W^2C_h)
(M_h^2\eta^{\mu\nu}-2k_2^\mu k_1^\nu)\\
&\hspace{18mm}-4M_W^2M_h^2\eta^{\mu\nu}
+2(M_h^2+6M_W^2)(2B_0-B_h)k_2^\mu k_1^\nu
\Bigr\},\\[2mm]
\int_p^{(d)}\mathscr M_3p^\mu k_1^\nu
={}&\frac{\mathcal N_W}{4M_W^2M_h^2}
\Bigl[
16M_W^2(M_h^2+6M_W^2)
-2(7M_h^2+48M_W^2)A_0\\
&\hspace{22mm}+(96M_W^4+2M_W^2M_h^2-M_h^4)B_h
\Bigr]k_2^\mu k_1^\nu,\\[2mm]
\int_p^{(d)}\mathscr M_4k_2^\mu p^\nu
={}&\frac{\mathcal N_W}{4M_W^2}
\Bigl[16M_W^2-18A_0+(M_h^2+14M_W^2)B_h\Bigr]k_2^\mu k_1^\nu,\\[2mm]
\int_p^{(d)}\mathscr M_5k_2^\mu k_1^\nu
={}&\mathcal N_W\Bigl[4B_0+16M_W^2C_h\Bigr]k_2^\mu k_1^\nu.
\end{align*}
这些式子展示了为什么不能在积分前把 $p^\mu p^\nu$ 直接替换为 $\eta^{\mu\nu}p^2/4$：各项分别含紫外发散的 $A_0$ 与 $B_0$，只有维数正规化下完成同一张量约化后，有限常数项才有唯一意义。

现在真正执行求和。相同质量的零动量二点函数满足
\begin{equation*}
 B_0(0;M_W^2,M_W^2)=\frac{A_0(M_W^2)}{M_W^2}-1.
\end{equation*}
把它代入上面五项，先收集 $\eta^{\mu\nu}$ 系数，再收集 $k_2^\mu k_1^\nu$ 系数。$k_2^\mu k_1^\nu$ 前的 $A_0$ 系数为
\begin{align*}
&\frac{4(M_h^2+6M_W^2)}{M_h^2M_W^2}
-\frac{7M_h^2+48M_W^2}{2M_h^2M_W^2}
-\frac{9}{2M_W^2}+\frac{4}{M_W^2}\\
&\qquad
=\frac{8(M_h^2+6M_W^2)-(7M_h^2+48M_W^2)-9M_h^2+8M_h^2}
{2M_h^2M_W^2}=0,
\end{align*}
而 $B_h$ 系数同样逐项给出
\begin{align*}
&-\frac{4(M_h^2+6M_W^2)}{M_h^2}
-\frac{M_h^4-2M_h^2M_W^2-96M_W^4}{4M_h^2M_W^2}
+\frac{M_h^2+14M_W^2}{4M_W^2}\\
&\qquad
=\frac{-16M_W^2(M_h^2+6M_W^2)
-(M_h^4-2M_h^2M_W^2-96M_W^4)
+M_h^2(M_h^2+14M_W^2)}{4M_h^2M_W^2}=0.
\end{align*}
因此所有一、二点母积分和紫外局域项都从物理振幅中消失。剩下的有限张量恰好组合为
\begin{equation*}
\mathcal M_W^{\mu\nu}
=\frac{\mathcal N_W}{M_h^2}
\Bigl[M_h^2+6M_W^2-6M_W^2(M_h^2-2M_W^2)C_h\Bigr]
\left(M_h^2\eta^{\mu\nu}-2k_2^\mu k_1^\nu\right).
\end{equation*}
由于
\begin{equation*}
M_h^2\eta^{\mu\nu}-2k_2^\mu k_1^\nu
=2\left(k_1\!\cdot k_2\,\eta^{\mu\nu}-k_2^\mu k_1^\nu\right),
\end{equation*}
两个光子 Ward 恒等式已经显式满足；横向性是全部图求和后的结果，而不是在单图上强加的投影。

最后只剩三点标量母积分。对三个等质量分母作 Feynman 参数化，正文\eqnrefs{eq:pv-scalar-masters-dp90} 给出
\begin{align*}
 C_h
&=-\int_0^1\dd x\int_0^{1-x}\dd y\,
\frac{1}{M_W^2-xyM_h^2-\ii0}\\
&=-\frac{2}{M_h^2}f(\tau_W),
\qquad
\tau_W\equiv\frac{4M_W^2}{M_h^2}=\frac{4m_W^2}{m_h^2}.
\end{align*}
代回有限系数并除去公共横向张量，得到
\begin{align*}
\frac{2}{M_h^2}
\Bigl[M_h^2+6M_W^2-6M_W^2(M_h^2-2M_W^2)C_h\Bigr]
&=2+3\tau_W+3\tau_W(2-\tau_W)f(\tau_W).
\end{align*}
按照正文对一圈函数的振幅符号约定，因此
\begin{equation*}
\boxed{
A_1(\tau)
=-\left[2+3\tau+3\tau(2-\tau)f(\tau)\right],}
\end{equation*}
即\eqnrefs{eq:sm-hgamma-loop-functions-dp41} 的第二行。用
$f(\tau)=\tau^{-1}+(3\tau^2)^{-1}+\order(\tau^{-3})$
还可立即检验
\begin{equation*}
A_1(\tau)=-7+\order(\tau^{-1}).
\end{equation*}
\end{derivation}

\begin{derivation}[从\eqnrefs{eq:qcd-singlet-axial-rotation-dp68}到\eqnrefs{eq:qcd-singlet-axial-anomaly-dp32}]
在欧几里得时空中取规范表示 $R$，并令
\begin{equation*}
 D_\mu=\partial_\mu+\ii g\mathcal A_\mu^aT^a,
 \qquad
 \Tr_R(T^aT^b)=T(R)\delta^{ab}.
\end{equation*}
用厄米算符 $\mathcal H\equiv\ii\slashed D_E$ 的正交本征函数展开 Grassmann 场：
\begin{equation*}
 \mathcal H\phi_n=\lambda_n\phi_n,
 \qquad
 \psi=\sum_na_n\phi_n,
 \qquad
 \bar\psi=\sum_n\bar b_n\phi_n^\dagger.
\end{equation*}
为了定义量子测度沿轴向方向的变化，取一参数局域轴向轨道
\begin{equation*}
 \psi_\tau=\ee^{\ii\tau\alpha(x)\gamma^5}\psi,
 \qquad
 \bar\psi_\tau=\bar\psi\,\ee^{\ii\tau\alpha(x)\gamma^5},
 \qquad \tau\in\mathbb R.
\end{equation*}
在本征函数系数空间中，$\psi$ 的变换矩阵为
\begin{equation*}
 C_{mn}(\tau)
 \equiv\int\dd^4x\,
 \phi_m^\dagger(x)\ee^{\ii\tau\alpha(x)\gamma^5}\phi_n(x),
 \qquad C(0)=I,
\end{equation*}
其在恒等元处的切向量严格为
\begin{equation*}
 \left.\frac{\dd C_{mn}}{\dd\tau}\right|_{\tau=0}
 =\ii\int\dd^4x\,
 \alpha(x)\phi_m^\dagger\gamma^5\phi_n.
\end{equation*}
Grassmann 测度按行列式的逆变换；$\psi$ 与 $\bar\psi$ 两份测度相乘。令 $J(\tau)$ 为总 Jacobian，则
\begin{align*}
 \left.\frac{\dd}{\dd\tau}\ln J(\tau)\right|_{\tau=0}
 &=-2\Tr\!\left[\left.\frac{\dd C}{\dd\tau}\right|_{0}\right]\\
 &=-2\ii\int\dd^4x\,\alpha(x)
 \sum_n\phi_n^\dagger(x)\gamma^5\phi_n(x).
\end{align*}
Ward 恒等式只需要这一沿群轨道的切向变化；随后对迹作规范协变正则化。
模和紫外发散，必须用保持规范协变性的同一算符谱正则化：
\begin{equation*}
 \mathscr D_R(x)
 \equiv\lim_{M\to\infty}
 \tr_{\rm spin,R}\!\left[
 \gamma^5\langle x|\ee^{-\mathcal H^2/M^2}|x\rangle
 \right].
\end{equation*}
把欧氏约定固定到每一个负号。令 $\mu=0,1,2,3$、
$\epsilon_E^{0123}=+1$，采用
\begin{align*}
\{\gamma_E^\mu,\gamma_E^\nu\}&=2\delta^{\mu\nu}I,\qquad
(\gamma_E^\mu)^\dagger=\gamma_E^\mu,\\
\gamma_E^5&=-\gamma_E^0\gamma_E^1\gamma_E^2\gamma_E^3,\qquad
\sigma_E^{\mu\nu}=\frac1{2\ii}[\gamma_E^\mu,\gamma_E^\nu].
\end{align*}
欧氏手征定向按后文的指数 $n_L-n_R$ 选择，即欧氏 $\gamma_E^5$ 的正本征空间
记为 $L$。它的正负标签须随 Wick 旋转一起匹配，不能直接代入闵氏
$P_L=(1-\gamma_M^5)/2$ 的投影符号。
欧氏 $\gamma$ 矩阵厄米，$D_\mu$ 与 $\slashed D_E$ 反厄米；
$M$ 是逆长度量纲的谱截断尺度。协变导数的曲率满足
$[D_\mu,D_\nu]=\ii g\mathcal F_{\mu\nu}^aT^a$，所以
\begin{align*}
\slashed D_E^2
&=D^2+\frac14[\gamma_E^\mu,\gamma_E^\nu][D_\mu,D_\nu]
=D^2-\frac g2\sigma_E^{\mu\nu}\mathcal F_{\mu\nu}^aT^a,\\
\mathcal H^2&=-(D^2+E),\qquad
E=-\frac g2\sigma_E^{\mu\nu}\mathcal F_{\mu\nu}^aT^a.
\end{align*}
Pauli 曲率项由 $\sigma_E^{\mu\nu}$ 的定义固定，手征迹则还依赖 $\gamma_E^5$ 的定向。若把 $\gamma_E^5$ 变号，手征迹及指数的左右手标签随之变号，而 Pauli 项不变。

对光滑规范背景，取无边界紧致调节，或在离边界固定距离的局域区域计算。
令 $t=M^{-2}$、$\Omega_{\mu\nu}=[D_\mu,D_\nu]$。Laplace 型算符的同点热核为
\begin{equation*}
\langle x|\ee^{-t\mathcal H^2}|x\rangle
=\frac1{(4\pi t)^2}\left[
I+tE+t^2\left(\frac12E^2+\frac16D^2E
+\frac1{12}\Omega_{\mu\nu}\Omega_{\mu\nu}\right)+\order(t^3)\right].
\end{equation*}
这里调用光滑背景的局域热核展开；其系数和适用边界条件参见
\href{https://arxiv.org/abs/hep-th/0306138}{Vassilevich 的热核展开综述}第 4、7 节。
该数学输入规定了曲率项、协变导数项与短时余项；以下直接对它们逐项取手征迹。
\begin{align*}
\tr_{\rm spin}(\gamma_E^5 I)&=0,&
\tr_{\rm spin}(\gamma_E^5 E)&=0,\\
\tr_{\rm spin}(\gamma_E^5D^2E)&=0,&
\tr_{\rm spin}(\gamma_E^5\Omega_{\mu\nu}\Omega_{\mu\nu})&=0.
\end{align*}
前两个等式因为不足四个不同的 $\gamma$ 矩阵；导数不会增加 $E$ 的自旋矩阵数；
$\Omega$ 则在自旋空间正比于单位矩阵。只有 $E^2/2$ 的迹留下四个不同指标。

由 $\sigma_E^{\mu\nu}=-\ii\gamma_E^\mu\gamma_E^\nu$（$\mu\ne\nu$）以及
$\gamma_E^5=-\gamma_E^0\gamma_E^1\gamma_E^2\gamma_E^3$，逐项反对易给
\begin{equation*}
\tr_{\rm spin}(\gamma_E^5\sigma_E^{\mu\nu}\sigma_E^{\rho\sigma})
=4\epsilon_E^{\mu\nu\rho\sigma}.
\end{equation*}
指标重复时迹为零；四指标不同时只由排列奇偶性决定符号。于是
\begin{align*}
\mathscr D_R(x)
&=\frac1{16\pi^2}\frac12\frac{g^2}{4}
\left(4\epsilon_E^{\mu\nu\rho\sigma}\right)
\Tr_R(T^aT^b)\mathcal F_{\mu\nu}^a\mathcal F_{\rho\sigma}^b\\
&=\frac{g^2T(R)}{32\pi^2}
\epsilon_E^{\mu\nu\rho\sigma}\mathcal F_{\mu\nu}^a\mathcal F_{\rho\sigma}^a\\
&=\frac{g^2T(R)}{16\pi^2}
\mathcal F_{\mu\nu}^a{}^\star\!\mathcal F_E^{a\mu\nu}.
\end{align*}
在固定光滑背景下，热核中余项乘上 $t^{-2}$ 后为 $\order(t)=\order(M^{-2})$，
所以去除截断的极限只保留上式的有限曲率二次项。
按上述手征与定向约定继续延拓到正文的闵氏指标，局域轴流 Ward 系数为
$g^2T(R)/(16\pi^2)$；场强、电场分量和四维积分测度必须沿同一 Wick 路径一起延拓。
\begin{equation*}
\mathscr D_R(x)
=\frac{g^2T(R)}{16\pi^2}
\mathcal F_{\mu\nu}^a{}^\star\!\mathcal F^{a\mu\nu}.
\end{equation*}
把这个极限代回测度的切向变化
\begin{equation*}
 \left.\frac{\dd}{\dd\tau}\ln J(\tau)\right|_{0}
 =-2\ii\int\dd^4x\,\alpha(x)\mathscr D_R(x),
\end{equation*}
再与质量项沿同一轴向轨道的经典变化相加，Ward 恒等式为
\begin{equation*}
 \partial_\mu J_5^\mu
 =\frac{2\ii mc}{\hbar}\bar\psi\gamma^5\psi
 +\frac{g^2\,2T(R)}{16\pi^2}
 \mathcal F_{\mu\nu}^a{}^\star\!\mathcal F^{a\mu\nu}.
\end{equation*}
QCD 基本表示满足 $T_F=1/2$；对 $N_f$ 个 Dirac 味求和即得到正文
\eqnrefs{eq:qcd-singlet-axial-anomaly-dp32}。两个 Grassmann 测度给出的因子 $2$
与基本表示的 $T_F=1/2$ 恰好相消。
对单位电荷阿贝尔表示则 $T(R)=1$，系数是 $g_{\rm em}^2/(8\pi^2)$。
再以 $F^{\rm geo}=\sqrt{\varepsilon_0c/\hbar}\,F$ 和
$g_{\rm em}=e/\sqrt{\varepsilon_0\hbar c}$ 恢复 SI 单位，
\begin{equation*}
\frac{g_{\rm em}^2}{8\pi^2}F^{\rm geo}{}^\star F^{\rm geo}
=\frac{e^2}{8\pi^2\hbar^2}F{}^\star F,
\end{equation*}
与正文\eqnrefs{eq:abj-anomaly-qed-dp11} 完全一致。反常来自有限的曲率二次热核系数；
光滑背景与规范协变调节是这一计算的假设，不能把它替换成不保持矢量规范对称的截断。
\end{derivation}

\begin{derivation}[从\eqnrefs{eq:electroweak-topological-charge-dp68}到\eqnrefs{eq:sm-bplusl-selection-dp32}]
先把拓扑信息与费米子路径积分连接起来。对一个处在 $SU(2)$ 表示 $R$ 的左手 Weyl 场，欧几里得 Dirac 算符的指数为
\begin{equation*}
 \operatorname{ind}\slashed D_E
 \equiv n_L-n_R=2T(R)\nu.
\end{equation*}
标准模型弱二重态属于 $R=\mathbf2$，$T(\mathbf2)=1/2$。在自对偶瞬子的半经典背景中，相反手征满足零模消失条件；当 $\nu>0$ 时，每个左手二重态因而有 $\nu$ 个左手零模而没有右手零模。一般光滑拓扑背景的指数只固定两者之差，允许额外的成对零模：
\begin{equation*}
 n_L-n_R=\nu.
\end{equation*}

把一个左手二重态沿 Dirac 本征函数展开，显式分离零模系数，
\begin{equation*}
 \psi_L(x)=\sum_{r=1}^{\nu}a_r\phi_{0r}(x)
 +\sum_{\lambda_n\neq0}a_n\phi_n(x).
\end{equation*}
费米子二次作用量只含非零本征值，因而测度中零模方向留下独立 Grassmann 积分
\begin{equation*}
 \int\mathcal D\psi_L\,\ee^{-S_F}
 =\mathcal N_{\rm nzm}
 \left(\prod_{r=1}^{\nu}\int\dd a_r\right),
 \qquad
 \mathcal N_{\rm nzm}\equiv
 \mathcal N_{\rm meas}\prod_{\lambda_n\neq0}\lambda_n,
\end{equation*}
Grassmann 规则 $\int\dd a_r=0$、$\int\dd a_r\,a_r=1$ 表明：没有外场插入时该拓扑扇区的费米子积分为零；一个非零相关函数必须为每个零模恰好提供一个相应的场因子。对 $\nu=1$，每个左手弱二重态必须出现一次。

一代标准模型含三个颜色副本 $Q_L^a$（$a=1,2,3$）和一个 $L_L$。把颜色收缩成单态、弱同位旋指标用 $\epsilon_{ij}$ 配对，零模饱和产生的局域多费米子结构可示意写成
\begin{equation*}
 \mathfrak O_{\rm tH}^{(g)}
 \equiv
 \epsilon_{abc}
 \bigl(Q_{Lg}^{a}Q_{Lg}^{b}\bigr)
 \bigl(Q_{Lg}^{c}L_{Lg}\bigr),
\end{equation*}
其中括号表示 Lorentz 与 $SU(2)_L$ 指标按不变张量收缩。三个夸克场各带 $B=1/3$，一个轻子场带 $L=1$，所以单代、单位拓扑荷的顶点改变
\begin{equation*}
 \Delta B_{g}=3\times\frac13=1,
 \qquad
 \Delta L_{g}=1.
\end{equation*}

同一结果也可直接从反常 Ward 恒等式积分得到。对重子流和轻子流，零模计数给每代相同的拓扑密度系数，
\begin{equation*}
 \partial_\mu j_B^\mu
 =N_g\frac{g^2}{32\pi^2}
 \mathcal W_{\mu\nu}^I{}^\star\!\mathcal W^{I\mu\nu},
 \qquad
 \partial_\mu j_L^\mu
 =N_g\frac{g^2}{32\pi^2}
 \mathcal W_{\mu\nu}^I{}^\star\!\mathcal W^{I\mu\nu}.
\end{equation*}
在一个连接两个渐近真空的四维区域上积分，并使用正文对 $\nu$ 的定义，
\begin{align*}
 \Delta B
 &=\int\dd^4x\,\partial_\mu j_B^\mu
 =N_g\nu,\\
 \Delta L
 &=\int\dd^4x\,\partial_\mu j_L^\mu
 =N_g\nu.
\end{align*}
对任意整数 $\nu$，$|\nu|$ 个零模副本给出同样的线性倍数；$\nu<0$ 时手征性与粒子数变化同时反向。于是
\begin{equation*}
 \Delta(B-L)=0,
 \qquad
 \Delta(B+L)=2N_g\nu,
\end{equation*}
即正文\eqnrefs{eq:sm-bplusl-selection-dp32}。选择定则来自费米子测度在不同拓扑扇区中的零模结构，而不是规范反常未抵消。
\end{derivation}

\begin{derivation}[从\eqnrefs{eq:qcd-singlet-axial-anomaly-dp32}到\eqnrefs{eq:strong-cp-theta-bar-dp11}]
把所有有色 Dirac 夸克收成味向量 $q=(u,d,s,\ldots)^T$，质量项写成
\begin{equation*}
 \Lag_m=-\bar q_LMq_R-\bar q_RM^\dagger q_L,
\end{equation*}
其中 $M$ 可以是任意非奇异复矩阵。取对角的实味空间矩阵
\begin{equation*}
 A={\rm diag}(\alpha_1,\ldots,\alpha_{N_f})
\end{equation*}
并作独立轴向重相位
\begin{equation*}
 q_L\longrightarrow \ee^{-\ii A}q_L,
 \qquad
 q_R\longrightarrow \ee^{+\ii A}q_R.
\end{equation*}
因为 $\bar q_L\to\bar q_Le^{+\ii A}$，质量矩阵在新场基底中变成
\begin{equation*}
 M' =\ee^{+\ii A}Me^{+\ii A}.
\end{equation*}
对行列式取相位，利用 $\det(\ee^{\ii A})=\ee^{\ii\Tr A}$：
\begin{align*}
 \det M'
 &=\ee^{\ii\Tr A}\det M\,\ee^{\ii\Tr A}
 =\ee^{2\ii\Tr A}\det M,\\
 \arg\det M'
 &=\arg\det M+2\Tr A
 \pmod{2\pi}.
\end{align*}
因此轴向旋转确实可以移动夸克质量矩阵的整体复相位；矢量型重相位 $q_L,q_R\to \ee^{\ii B}q_{L,R}$ 则不会产生这一相位变化。

另一方面，正文\eqnrefs{eq:qcd-singlet-axial-anomaly-dp32} 表明同一个轴向变换不是量子测度的对称。正文\eqnrefs{eq:qcd-theta-action-dp89} 已固定 $\theta_s$ 的归一化。对第 $i$ 个基本表示 Dirac 味作角度 $\alpha_i$ 的轴向旋转，Fujikawa Jacobian 给
\begin{equation*}
 \ln J_i
 =-\ii\alpha_i\frac{g_s^2}{16\pi^2}
 \int\dd^4x\,
 \mathcal F_{\mu\nu}^a{}^\star\!\mathcal F^{a\mu\nu}.
\end{equation*}
把 $J_i$ 吸收到 $\exp(\ii S/\hbar)$ 中，等价于
\begin{equation*}
 \theta_s\longrightarrow\theta_s-2\alpha_i.
\end{equation*}
全部味同时变换时便有
\begin{equation*}
 \theta_s' =\theta_s-2\Tr A.
\end{equation*}
与质量行列式的变化相加：
\begin{align*}
 \theta_s'+\arg\det M'
 &=(\theta_s-2\Tr A)
 +(\arg\det M+2\Tr A)\\
 &=\theta_s+\arg\det M.
\end{align*}
所以轴向基底选择能分别改变 $\theta_s$ 与质量矩阵的整体相位，却不能改变组合
\begin{equation*}
 \boxed{\bar\theta=\theta_s+\arg\det M}.
\end{equation*}
对标准模型六个夸克，把 $M$ 分成上、下型两个块就得到正文\eqnrefs{eq:strong-cp-theta-bar-dp11} 的 $\arg\det(M_uM_d)$。

若某个夸克严格无质量，例如某个质量本征值 $m_j=0$，则相应轴向角 $\alpha_j$ 不会在质量项中产生可观测相位：$\det M=0$ 时整体行列式相位本身不再是物理参数。此时可以选择 $\alpha_j=\theta_s/2$ 把拓扑角完全移入这个无质量场的轴向定义，从而使 $\theta_s$ 不可观测。真实 QCD 中所有夸克质量均非零，因此这一消除机制不存在，$\bar\theta$ 成为独立的强 $CP$ 参数。
\end{derivation}



\begin{derivation}[从\eqnrefs{eq:eeww-three-amplitudes-dp11}到\eqnrefs{eq:eeww-right-cancellation-dp11}]
令中性传播子的入射四动量为 $q=p_1+p_2=k_-+k_+$。对三个动量都取向顶角流入的 Yang--Mills 张量，散射图中需要
\begin{equation*}
\Gamma_{\rho\mu\nu}(-q,k_-,k_+)
=(k_--k_+)_\rho g_{\mu\nu}
 +(k_++q)_\mu g_{\nu\rho}
 -(q+k_-)_\nu g_{\rho\mu}.
\end{equation*}
两个末态 $W$ 在壳，$k_-^2=k_+^2=m_W^2c^2$，并且
\begin{equation*}
q^2=s=2m_W^2c^2+2k_-\!\cdot k_+.
\end{equation*}
由精确分解~\eqnrefs{eq:longitudinal-polarization-highE-dp11}，完整双纵向收缩恒等地分成 $kk$、$kr_L$、$r_Lk$ 与 $r_Lr_L$ 四部分。先提取 $kk$ 部分，即在两个外腿上分别取 $k_-^\mu/(m_Wc)$ 与 $k_+^\nu/(m_Wc)$；顶角的三个张量项分别给
\begin{align*}
\Gamma_{\rho\mu\nu}k_-^\mu k_+^\nu
={}&(k_-\!\cdot k_+)(k_--k_+)_\rho\\
&+\bigl[k_-\!\cdot(k_++q)\bigr]k_{+\rho}
-\bigl[k_+\!\cdot(q+k_-)\bigr]k_{-\rho}.
\end{align*}
利用 $q=k_-+k_+$，两方括号都是 $m_W^2c^2+2k_-\!\cdot k_+$，于是
\begin{align*}
\Gamma_{\rho\mu\nu}k_-^\mu k_+^\nu
&=\bigl(m_W^2c^2+k_-\!\cdot k_+\bigr)(k_+-k_-)_\rho\\
&=\frac{s}{2}(k_+-k_-)_\rho .
\end{align*}
无质量电子的手征流满足
\begin{equation*}
q_\rho J_\chi^\rho
=\bar v(p_2)(\slashed p_1+\slashed p_2)P_\chi u(p_1)=0,
\end{equation*}
其中分别用了 $\slashed p_1u(p_1)=0$ 与 $\bar v(p_2)\slashed p_2=0$。因此 $J_\chi\!\cdot k_-=-J_\chi\!\cdot k_+$，从而
\begin{equation*}
\frac{1}{s}J_\chi^\rho
\Gamma_{\rho\mu\nu}k_-^\mu k_+^\nu
=J_\chi\!\cdot k_+.
\end{equation*}
光子图与 $Z$ 图的双 $k/(m_Wc)$ 项具有同一个运动学因子 $\mathcal K_\chi$。所有至少含一个 $r_L$ 的项由~\eqnrefs{eq:longitudinal-polarization-highE-dp11} 的范数界控制，在固定角高能极限中消失。

对右手电子，$P_Lu_R=0$，所以 $t$ 道中微子图严格消失。由
\begin{equation*}
g_{WW\gamma}=gs_W=g_{\rm em},\qquad
 g_{WWZ}=gc_W,\qquad
 g_{\gamma e_R}=-g_{\rm em},\qquad
 g_{Ze_R}=\frac{g}{c_W}s_W^2,
\end{equation*}
两个 $s$ 道主导项之和为
\begin{align*}
\mathcal M_{\gamma,R}^{(E^2)}+\mathcal M_{Z,R}^{(E^2)}
&=\left[g_{WW\gamma}g_{\gamma e_R}+g_{WWZ}g_{Ze_R}\right]\mathcal K_R\\
&=\left[-g^2s_W^2+g^2s_W^2\right]\mathcal K_R=0,
\end{align*}
即正文\eqnrefs{eq:eeww-right-cancellation-dp11}。这里的相消同时使用了 $e=gs_W$、$WWZ$ 的非阿贝尔顶角系数和右手电子的弱同位旋 $T^3=0$；任意改变其中一项都会留下随能量增长的非物理振幅。
\end{derivation}

\begin{derivation}[从\eqnrefs{eq:eeww-three-amplitudes-dp11}到\eqnrefs{eq:eeww-left-total-cancellation-dp11}]

左手中性流耦合为
\begin{equation*}
g_{Ze_L}=\frac g{c_W}\left(-\frac12+s_W^2\right).
\end{equation*}
由正文\eqnrefs{eq:eeww-common-kinematic-structure-dp75} 前的 Yang--Mills 双纵向收缩，两个 $s$ 道图共享 $\mathcal K_L$，故
\begin{align*}
\mathcal M_{L,s}^{(E^2)}
&=\left[g_{WW\gamma}(-g_{\rm em})+g_{WWZ}g_{Ze_L}\right]\mathcal K_L\\
&=\left[-g^2s_W^2+g^2\left(-\frac12+s_W^2\right)\right]\mathcal K_L\\
&=-\frac{g^2}{2}\mathcal K_L.
\end{align*}

再看中微子图。取 $r=p_1-k_-$，其张量部分为
\begin{equation*}
\mathcal M_\nu^{\mu\nu}
=-\frac{g^2}{2}\bar v(p_2)\gamma^\nu P_L
\frac{\slashed r}{r^2}\gamma^\mu P_Lu(p_1).
\end{equation*}
同样提取中微子图的双 $k/(m_Wc)$ 系数，分别取 $k_-^\mu/(m_Wc)$ 与 $k_+^\nu/(m_Wc)$。由 $\slashed p_1u(p_1)=0$，
\begin{equation*}
\slashed k_-P_Lu=(\slashed p_1-\slashed r)P_Lu
=-\slashed rP_Lu.
\end{equation*}
因此传播子分子与右侧收缩给出的 $\slashed r$ 合并，利用 $\slashed r\slashed r=r^2$ 消去传播子分母：
\begin{align*}
\mathcal M_{L,t}^{(E^2)}
&=\frac{g^2}{2m_W^2c^2}
\bar v(p_2)\slashed k_+P_Lu(p_1)\\
&=\frac{g^2}{2}\mathcal K_L.
\end{align*}
$\mathcal K_L$ 正是两个 $s$ 道已经出现的同一电子流结构。于是双动量系数满足
\begin{equation*}
\mathcal M_{L,s}^{(E^2)}+\mathcal M_{L,t}^{(E^2)}
=\left(-\frac{g^2}{2}+\frac{g^2}{2}\right)\mathcal K_L=0,
\end{equation*}
得到正文\eqnrefs{eq:eeww-left-total-cancellation-dp11}。
\end{derivation}

\begin{derivation}[从\eqnrefs{eq:sm-hgg-amplitude-dp94}到\eqnrefs{eq:sm-hgg-effective-operator-dp94}]
先对单个夸克味 $q$ 计算完整三角圈，而不预先代入重顶极限。Higgs--夸克约化顶角为 $-\ii m_qc^2/v_E$，夸克--胶子顶角为 $-\ii g_s\gamma^\mu T^a$。取两个出射胶子动量为 $k_1,k_2$，满足
\begin{equation*}
k_1^2=k_2^2=0,
\qquad
2k_1\!\cdot k_2=m_h^2c^2.
\end{equation*}
两种外胶子排序都必须保留。提出共同外腿归一化后，第一种排序的圈张量可写成
\begin{align*}
\mathcal M_{q,1}^{ab,\mu\nu}
={}&-g_s^2\frac{m_qc^2}{v_E}
\Tr(T^aT^b)
\int_\ell^{(d)}
\frac{\Tr\!\left[
(\slashed\ell+m_qc)\gamma^\mu
(\slashed\ell+\slashed k_1+m_qc)\gamma^\nu
(\slashed\ell-\slashed k_2+m_qc)
\right]}
{D_0D_1D_2},
\end{align*}
其中
\begin{align*}
D_0&=\ell^2-m_q^2c^2+\ii0,\\
D_1&=(\ell+k_1)^2-m_q^2c^2+\ii0,\\
D_2&=(\ell-k_2)^2-m_q^2c^2+\ii0.
\end{align*}
第二种排序由 $(a,\mu,k_1)\leftrightarrow(b,\nu,k_2)$ 得到。颜色部分立即给
\begin{equation*}
\Tr(T^aT^b)+\Tr(T^bT^a)
=2T_F\delta^{ab}=\delta^{ab},
\end{equation*}
所以一个颜色单态 Higgs 只能产生 $\delta^{ab}$ 颜色结构。

现在真正做圈积分。三个分母用 Feynman 参数合并：
\begin{equation*}
\frac1{D_0D_1D_2}
=2\int_0^1\dd x\int_0^{1-x}\dd y\,
\frac1{[xD_1+yD_2+(1-x-y)D_0]^3}.
\end{equation*}
令
\begin{equation*}
L=\ell+xk_1-yk_2,
\end{equation*}
则
\begin{equation*}
xD_1+yD_2+(1-x-y)D_0
=L^2-\Delta_q(x,y)+\ii0,
\end{equation*}
其中
\begin{equation*}
\Delta_q(x,y)=m_q^2c^2-xy\,m_h^2c^2.
\end{equation*}
因此阈值结构已经显式出现：当参数区域内 $\Delta_q=0$ 时，对应 $m_h=2m_q$ 的实夸克对分支点。

分子不能只保留质量项。记 $M_q=m_qc$ 以及
\begin{equation*}
a=\ell,\qquad b=\ell+k_1,\qquad c=\ell-k_2.
\end{equation*}
第一张图的 Dirac 迹为
\begin{align*}
N^{\mu\nu}
&=\Tr[(\slashed a+M_q)\gamma^\mu
(\slashed b+M_q)\gamma^\nu(\slashed c+M_q)].
\end{align*}
奇数个 $\gamma$ 的项为零，只剩一个 $M_q$ 与三个 $M_q$ 的项。逐项使用四 $\gamma$ 迹公式可得
\begin{align*}
\frac{N^{\mu\nu}}{4M_q}
={}&b^\mu c^\nu+c^\mu b^\nu-\eta^{\mu\nu}b\!\cdot c
+a^\mu c^\nu-a^\nu c^\mu+\eta^{\mu\nu}a\!\cdot c\\
&+a^\mu b^\nu+a^\nu b^\mu-\eta^{\mu\nu}a\!\cdot b
+M_q^2\eta^{\mu\nu}.
\end{align*}
平移到 $L$ 后，奇数次 $L$ 积分消失，秩二积分按
\begin{equation*}
\int\dd^dL\,L^\mu L^\nu F(L^2)
=\frac{\eta^{\mu\nu}}{d}\int\dd^dL\,L^2F(L^2)
\end{equation*}
约化。两张排序图相加时，非横向的局域二点型部分精确相消；剩余张量必满足
\begin{equation*}
k_1^\mu\mathcal M_{\mu\nu}^{ab}=0,
\qquad
k_2^\nu\mathcal M_{\mu\nu}^{ab}=0,
\end{equation*}
所以可以写成
\begin{equation*}
\mathcal M_{q}^{ab,\mu\nu}
=\delta^{ab}C_q
\left(k_1\!\cdot k_2\,\eta^{\mu\nu}-k_2^\mu k_1^\nu\right).
\end{equation*}
把显式分子投影到这个横向张量后，所有圈动量积分化为一个参数积分：
\begin{equation*}
A_{1/2}(\tau_q)
=4\tau_q\int_0^1\dd x\int_0^{1-x}\dd y\,
\frac{1-4xy}{\tau_q-4xy-\ii0},
\qquad
\tau_q\equiv\frac{4m_q^2}{m_h^2}.
\end{equation*}
这一步与前文 $h\to\gamma\gamma$ 的费米子环具有相同 Lorentz 积分，但这里的电荷因子 $N_cQ_q^2g_{\rm em}^2$ 被 QCD 颜色迹与 $g_s^2$ 取代。令
\begin{equation*}
I(\tau)=\int_0^1\dd x\int_0^{1-x}\dd y\,
\frac1{\tau-4xy-\ii0},
\end{equation*}
先积掉 $y$：
\begin{equation*}
I(\tau)=\int_0^1\frac{\dd x}{4x}
\ln\!\left[\frac{\tau}{\tau-4x(1-x)-\ii0}\right].
\end{equation*}
定义
\begin{equation*}
f(\tau)=2I(\tau)
=\begin{cases}
\arcsin^2(1/\sqrt\tau),&\tau\ge1,\\[4pt]
-\dfrac14\left[
\ln\dfrac{1+\sqrt{1-\tau}}{1-\sqrt{1-\tau}}-\ii\pi
\right]^2,&\tau<1,
\end{cases}
\end{equation*}
再把 $1-4xy=(\tau-4xy)+(1-\tau)$ 分开，得到
\begin{align*}
A_{1/2}(\tau)
&=4\tau\left[\frac12+(1-\tau)I(\tau)\right]\\
&=2\tau\left[1+(1-\tau)f(\tau)\right].
\end{align*}
因此单味系数为
\begin{equation*}
C_q=\frac{\alpha_s}{4\pi v_E}A_{1/2}(\tau_q),
\end{equation*}
对所有夸克味求和即得到正文\eqnrefs{eq:sm-hgg-amplitude-dp94}。这里的 $A_{1/2}$ 不是借用的经验函数，而是三角圈参数积分的闭式。

由振幅进入衰变宽度时，记
\begin{equation*}
C_{gg}=\frac{\alpha_s}{4\pi v_E}\sum_qA_{1/2}(\tau_q),
\qquad
T_{\mu\nu}=k_1\!\cdot k_2\,\eta_{\mu\nu}-k_{2\mu}k_{1\nu}.
\end{equation*}
Ward 恒等式允许用 $-\eta_{\mu\nu}$ 完成物理偏振求和。颜色求和与 Lorentz 收缩分别给
\begin{equation*}
\sum_{a,b}\delta^{ab}\delta^{ab}=N_c^2-1=8,
\qquad
T_{\mu\nu}T^{\mu\nu}=2(k_1\!\cdot k_2)^2
=\frac{m_h^4c^4}{2}.
\end{equation*}
故
\begin{equation*}
\sum|\mathcal M|^2=4|C_{gg}|^2m_h^4c^4.
\end{equation*}
两个末态胶子完全相同，二体相空间必须除以 $2!$：
\begin{equation*}
\Gamma
=\frac1{2!}\frac1{2m_hc^2}\frac1\hbar
\int\dd\Phi_2\sum|\mathcal M|^2,
\qquad
\int\dd\Phi_2=\frac1{8\pi}.
\end{equation*}
再用 $G_F^{(E)}/\sqrt2=1/(2v_E^2)$，便得到正文\eqnrefs{eq:sm-hgg-width-dp94}。

最后才取重顶受控极限。$\tau_t\gg1$ 时
\begin{equation*}
f(\tau_t)=\frac1{\tau_t}+\frac1{3\tau_t^2}+\order(\tau_t^{-3}),
\end{equation*}
所以
\begin{equation*}
A_{1/2}(\tau_t)=\frac43+\frac{14}{45\tau_t}+\order(\tau_t^{-2}).
\end{equation*}
若写局域理论
\begin{equation*}
\Lag_{\rm eff}=C_h\frac{h}{v_E}G^a_{\mu\nu}G^{a\mu\nu},
\end{equation*}
其两胶子树级矩阵元为
\begin{equation*}
\mathcal M_{\mu\nu}^{ab}
=4C_h\frac{\delta^{ab}}{v_E}
\left(k_1\!\cdot k_2\,\eta_{\mu\nu}-k_{2\mu}k_{1\nu}\right).
\end{equation*}
与完整一圈结果的首项逐项匹配：
\begin{equation*}
4C_h\frac1{v_E}
=\frac{\alpha_s}{4\pi v_E}\frac43,
\end{equation*}
从而
\begin{equation*}
\boxed{C_h=\frac{\alpha_s}{12\pi}.}
\end{equation*}
下一项相对修正由 $1/\tau_t=m_h^2/(4m_t^2)$ 控制，所以局域算符不是“把顶夸克删掉”的猜测，而是完整三角圈在低外动量下的系统展开。
\end{derivation}

\begin{derivation}[从带电流平均场到物质势]
为跟踪量纲，取规范归一化空间场 $\psi$ 的量纲为 $\mathrm{m^{-3/2}}$。低能带电流经 Fierz 重排后，其哈密顿密度为
\begin{equation*}
 \mathcal H_{\rm CC}^{\rm eff}
 =\frac{G_F^{(E)}(\hbar c)^3}{\sqrt2}
 [\bar\nu_e\gamma^\mu(1-\gamma^5)\nu_e]
 [\bar e\gamma_\mu(1-\gamma^5)e].
\end{equation*}
系数的单位是能量乘体积，两个流的乘积单位是 $\mathrm{m^{-6}}$，所以总量为能量密度。介质静止、无极化且忽略正电子时，
\begin{equation*}
 \langle\bar e\gamma^\mu e\rangle=(n_e,\boldsymbol0),\qquad
 \langle\bar e\gamma^\mu\gamma^5e\rangle=0.
\end{equation*}
对传播的左手中微子，$1-\gamma^5=2P_L$，因而
\begin{equation*}
 \langle\mathcal H_{\rm CC}^{\rm eff}\rangle_e
 =\sqrt2G_F^{(E)}(\hbar c)^3n_e\,\nu_{eL}^\dagger\nu_{eL}.
\end{equation*}
这就确定了\eqnrefs{eq:sm-matter-flavor-hamiltonian}的势及正号。有正电子时树级矢量流数密度改为净密度 $n_{e^-}-n_{e^+}$；反中微子的相对势反号。

味基与质量基的方向也可直接检查。设 $\ell_L^{\rm weak}=U_{eL}\ell_L$、$\nu_L^{\rm weak}=U_\nu\nu_L$，则
\begin{equation*}
 \bar\nu_L^{\rm weak}\gamma^\mu\ell_L^{\rm weak}
 =\bar\nu_{iL}\gamma^\mu(U_\nu^\dagger U_{eL})_{i\alpha}\ell_{\alpha L}
 =\bar\nu_{iL}\gamma^\mu U_{\alpha i}^*\ell_{\alpha L}.
\end{equation*}
故产生味态的质量振幅是 $U_{\alpha i}^*$，探测为 $\beta$ 味的投影是 $U_{\beta i}$；组合成 $U\operatorname{diag}(m_i^2)U^\dagger$ 与真空振幅的顺序一致。
\end{derivation}







\begin{derivation}[从\eqnrefs{eq:general-gauge-beta-weyl-scalar-dp11}到\eqnrefs{eq:sm-gauge-beta-one-loop-dp11}]
正文\eqnrefs{eq:general-gauge-beta-weyl-scalar-dp11} 已给出一般规范场、Weyl 费米子和复标量对一圈 $\beta$ 系数的群论权重。把标准模型三代手征场与一个 Higgs 二重态逐项代入，并正确乘上颜色和弱同位旋多重度，就得到正文\eqnrefs{eq:sm-gauge-beta-one-loop-dp11}。

对非阿贝尔规范群，一圈系数采用
\begin{equation*}
b=-\frac{11}{3}C_A
+\frac23\sum_{\text{左手 Weyl 费米子}}T(R_f)
+\frac13\sum_{\text{复标量}}T(R_s).
\end{equation*}
每个求和都必须包含其他内部量子数产生的多重度。

对 $SU(3)_c$，$C_A=3$。每一代中 $Q_L$ 含两个颜色基本表示，$u_R$ 与 $d_R$ 各含一个，因此三代共有 $12$ 个 Weyl 基本表示。用 $T(\mathbf3)=1/2$，
\begin{align*}
\sum_fT_3(R_f)&=12\times\frac12=6,\\
b_3&=-\frac{11}{3}(3)+\frac23(6)\\
&=-7.
\end{align*}
Higgs 是颜色单态，不贡献 $SU(3)_c$ 项。

对 $SU(2)_L$，$C_A=2$。每一代有三个颜色副本的 $Q_L$ 双重态和一个 $L_L$ 双重态，三代共 $12$ 个 Weyl 双重态，因此
\begin{equation*}
\sum_fT_2(R_f)=12\times\frac12=6.
\end{equation*}
Higgs 是一个复双重态，$T_2(\mathbf2)=1/2$。于是
\begin{align*}
b_2
&=-\frac{11}{3}(2)+\frac23(6)+\frac13\left(\frac12\right)\\
&=-\frac{19}{6}.
\end{align*}

对 $U(1)_Y$ 没有规范玻色子自相互作用项；表示指标改为所有内部多重度加权的 $Y^2$。每一代费米子给出
\begin{align*}
\sum_fd_3d_2Y^2
={}&6\left(\frac16\right)^2
+3\left(\frac23\right)^2
+3\left(-\frac13\right)^2\\
&+2\left(-\frac12\right)^2+(-1)^2\\
={}&\frac{10}{3}.
\end{align*}
三代因此给 $10$。Higgs 的两个复分量均有 $Y=1/2$，所以标量和为 $2(1/2)^2=1/2$。于是
\begin{align*}
b_Y
&=\frac23(10)+\frac13\left(\frac12\right)\\
&=\frac{41}{6}.
\end{align*}
三组结果合并即得到\eqnrefs{eq:sm-gauge-beta-one-loop-dp11}。
\end{derivation}



\begin{derivation}[从\eqnrefs{eq:sm-gim-loop-sum-dp83}到\eqnrefs{eq:sm-gim-mass-suppression-dp83}]
质量本征基中的夸克带电流为
\begin{equation*}
 \Lag_{\rm CC}^{q}
 =-\frac{\hbar c\,g}{\sqrt2}
 \left[
 \bar u_i\gamma^\mu P_LV_{ij}d_j\,\mathcal W_\mu^+
 +\bar d_j\gamma^\mu P_LV_{ij}^*u_i\,\mathcal W_\mu^-
 \right].
\end{equation*}
考虑 $d_j\to d_k$、$j\neq k$ 的一圈味改变中性流。Dyson 展开中任意包含两个带电流顶角的收缩，都含味因子
\begin{equation*}
 V_{ij}V_{ik}^*=V_{ik}^*V_{ij},
\end{equation*}
因为第一个顶角把外部 $d_j$ 连接到内部 $u_i$，第二个厄米共轭顶角再把同一内部味 $u_i$ 连接到外部 $d_k$。

给定外部粒子与动量后，选一组完备、线性无关的 Lorentz--Dirac 算符基 $\{\mathcal K_a(p_{\rm ext})\}$。同一微扰阶的全部规范完成图求和后，任意内部味 $i$ 的贡献都能按这组基展开：
\begin{equation*}
 \mathcal A_{j\to k}^{(i)}
 =V_{ik}^*V_{ij}
 \sum_a\mathcal K_a(p_{\rm ext})F_a(x_i),
 \qquad
 x_i=\frac{m_{u_i}^2}{m_W^2}.
\end{equation*}
这里 $F_a$ 由实际圈积分定义。例如在维数正规化中，每个 $F_a$ 都来自形如
\begin{equation*}
 \mu_{\rm DR}^{2\epsilon}
 \int\frac{\dd^d\ell}{(2\pi)^d}
 \frac{N_a(\ell,p_{\rm ext};m_{u_i},m_W)}
 {\prod_r D_r(\ell,p_{\rm ext};m_{u_i},m_W)}
\end{equation*}
的积分及同阶规范完成项。对内部味求和并交换有限求和与算符基展开，得到
\begin{align*}
 \mathcal A_{j\to k}^{\rm FCNC}
 &=\sum_i\mathcal A_{j\to k}^{(i)}\\
 &=\sum_a\mathcal K_a(p_{\rm ext})
 \left[\sum_iV_{ik}^*V_{ij}F_a(x_i)\right].
\end{align*}
因此
\begin{equation*}
 C_a^{jk}=\sum_iV_{ik}^*V_{ij}F_a(x_i),
\end{equation*}
即正文\eqnrefs{eq:sm-gim-loop-sum-dp83}。

CKM 幺正性 $V^\dagger V=I$ 的非对角元满足
\begin{equation*}
 \sum_iV_{ik}^*V_{ij}=0,
 \qquad j\neq k.
\end{equation*}
对任意参考点 $x_*$，逐个算符系数加减同一个 $F_a(x_*)$：
\begin{align*}
 C_a^{jk}
 &=\sum_iV_{ik}^*V_{ij}
 \left[F_a(x_i)-F_a(x_*)+F_a(x_*)\right]\\
 &=\sum_iV_{ik}^*V_{ij}[F_a(x_i)-F_a(x_*)]
 +F_a(x_*)\sum_iV_{ik}^*V_{ij}\\
 &=\sum_iV_{ik}^*V_{ij}[F_a(x_i)-F_a(x_*)].
\end{align*}
这就是\eqnrefs{eq:sm-gim-subtracted-exact-dp85}。若 $x_i=x_*$ 对所有内部味都成立，则每个独立算符系数 $C_a^{jk}$ 都等于零，因而完整振幅为零。

若 $F_a$ 在连接 $x_*$ 与 $x_i$ 的区间可微，微积分基本定理还给出精确质量差表示
\begin{align*}
 F_a(x_i)-F_a(x_*)
 &=\int_{x_*}^{x_i}\dd y\,F_a'(y)\\
 &=(x_i-x_*)\int_0^1\dd t\,
 F_a'\!\left[x_*+t(x_i-x_*)\right],
\end{align*}
因此
\begin{equation*}
 C_a^{jk}
 =\sum_iV_{ik}^*V_{ij}(x_i-x_*)
 \int_0^1\dd t\,
 F_a'\!\left[x_*+t(x_i-x_*)\right].
\end{equation*}
这条等式直接表明共同的内部质量部分被消去，剩余量由质量劈裂加权。

对轻内部质量取 $x_*=0$。若 $F_a\in C^3([0,x_{\max}])$，Taylor 定理带积分余项给
\begin{equation*}
 F_a(x_i)
 =F_a(0)+x_iF_a'(0)+\frac{x_i^2}{2}F_a''(0)
 +\frac{x_i^3}{2}\int_0^1\dd t\,(1-t)^2F_a^{(3)}(tx_i).
\end{equation*}
代入并使用 $\sum_iV_{ik}^*V_{ij}=0$：
\begin{equation*}
 C_a^{jk}
 =F_a'(0)\sum_iV_{ik}^*V_{ij}x_i
 +\mathcal R_{a,jk}^{(\ge2)},
\end{equation*}
其中
\begin{align*}
 \mathcal R_{a,jk}^{(\ge2)}
 &\equiv
 \frac{F_a''(0)}2
 \sum_iV_{ik}^*V_{ij}x_i^2+R_{a,jk}^{(3)},\\
 R_{a,jk}^{(3)}
 &\equiv
 \frac12\sum_iV_{ik}^*V_{ij}x_i^3
 \int_0^1\dd t\,(1-t)^2F_a^{(3)}(tx_i).
\end{align*}
于是
\begin{align*}
 |R_{a,jk}^{(3)}|
 &\le\frac16
 \max_{0\le y\le x_{\max}}|F_a^{(3)}(y)|
 \sum_i|V_{ik}^*V_{ij}|x_i^3,\\
 |\mathcal R_{a,jk}^{(\ge2)}|
 &\le\frac{|F_a''(0)|}{2}
 \sum_i|V_{ik}^*V_{ij}|x_i^2
 +|R_{a,jk}^{(3)}|.
\end{align*}
最后用 $x_i=m_{u_i}^2/m_W^2$ 即得到正文\eqnrefs{eq:sm-gim-mass-suppression-dp83} 及其余项控制。若 $F_a$ 在零点附近非解析，或某个 $x_i$ 不小，则保留精确差值\eqnrefs{eq:sm-gim-subtracted-exact-dp85}。
\end{derivation}

\begin{derivation}[从\eqnrefs{eq:sm-full-lagrangian-blocks-dp11}到\eqnrefs{eq:minimal-sm-parameter-count-dp11}]
Yukawa 矩阵的矩阵元依赖味基选择，参数计数必须先除去这种场重定义冗余，因此从夸克扇区开始。三代时，$Y_u$ 与 $Y_d$ 各是一般复 $3\times3$ 矩阵。每个矩阵含

\begin{equation*}
 9\ \text{个复矩阵元}
 =18\ \text{个实参数},
\end{equation*}
所以两者合计有 $36$ 个实参数。

暂时令 Yukawa 耦合为零。夸克动能项在世代空间具有
\begin{equation*}
 U(3)_{Q_L}\times U(3)_{u_R}\times U(3)_{d_R}
\end{equation*}
的味基冗余。每个 $U(3)$ 有 $3^2=9$ 个连续生成元，因此三组基变换一共有 $27$ 个连续参数。不过不能直接用 $36-27$，因为一般非零的 $Y_u,Y_d$ 仍保留共同的矢量式重子数变换
\begin{equation*}
 Q_L\to \ee^{\ii\alpha}Q_L,
 \qquad
 u_R\to \ee^{\ii\alpha}u_R,
 \qquad
 d_R\to \ee^{\ii\alpha}d_R.
\end{equation*}
在 Yukawa 双线性项中，左手场经 Dirac 共轭后产生 $\ee^{-\ii\alpha}$，恰与右手场的 $\ee^{+\ii\alpha}$ 相消。因此 $U(1)_B$ 是打开 Yukawa 耦合后仍未破缺的子群，不能用来改变 $Y_u,Y_d$。真正能够消去坐标冗余的味基参数只有
\begin{equation*}
 27-1=26.
\end{equation*}
所以夸克 Yukawa 扇区留下
\begin{equation*}
 36-26=10
\end{equation*}
个物理参数。为了看清这 $10$ 个参数为什么恰好表现为“六个质量、三个角、一个相位”，还要把上述商空间计数落实到质量本征基。对两个一般复矩阵分别作奇异值分解：
\begin{equation*}
 Y_u=U_{uL}\,y_u\,U_{uR}^\dagger,
 \qquad
 Y_d=U_{dL}\,y_d\,U_{dR}^\dagger,
\end{equation*}
其中
\begin{equation*}
 y_u=\operatorname{diag}(y_u,y_c,y_t),
 \qquad
 y_d=\operatorname{diag}(y_d,y_s,y_b)
\end{equation*}
可取为非负实对角矩阵。右手幺正矩阵 $U_{uR},U_{dR}$ 把两个质量矩阵对角化；带电弱流只含左手双重态，因此两个左手旋转之间的相对矩阵保留下来：
\begin{equation*}
 V_{\rm CKM}=U_{uL}^\dagger U_{dL}.
\end{equation*}
六个奇异值与六个夸克质量一一对应，剩下的问题就是一个一般 $3\times3$ 幺正矩阵究竟含多少可观测参数。

一般 $U(3)$ 有 $9$ 个实参数。可把它们组织成三个实混合角与六个相位。质量项已经对角以后，六个 Dirac 夸克仍可分别作共同左右手重相位
\begin{equation*}
 u_i\to \ee^{\ii\alpha_i}u_i,
 \qquad
 d_j\to \ee^{\ii\beta_j}d_j,
\end{equation*}
而不改变任何质量本征值。带电流
$\bar u_{Li}\gamma^\mu V_{ij}d_{Lj}$ 因而只把 CKM 元素变换为
\begin{equation*}
 V_{ij}\longrightarrow
 \ee^{-\ii\alpha_i}V_{ij}\ee^{+\ii\beta_j}.
\end{equation*}
六个相位 $\{\alpha_i,\beta_j\}$ 中，若全部同时增加同一个 $\omega$，则
\begin{equation*}
 \ee^{-\ii(\alpha_i+\omega)}V_{ij}
 \ee^{+\ii(\beta_j+\omega)}
 =\ee^{-\ii\alpha_i}V_{ij}\ee^{+\ii\beta_j},
\end{equation*}
所以共同的重子数相位并不能改变 $V_{\rm CKM}$。真正可用于消去 CKM 相位的只有 $6-1=5$ 个独立重相位。因此六个原始相位中只剩
\begin{equation*}
 6-5=1
\end{equation*}
个物理 $CP$ 相位，而三个实混合角保留。于是夸克 Yukawa 扇区的 $10$ 个参数被显式分解为
\begin{equation*}
 \boxed{6\ \text{个夸克质量}
 +3\ \text{个 CKM 混合角}
 +1\ \text{个 CKM 相位}.}
\end{equation*}

带电轻子扇区只有一个一般复 $3\times3$ 矩阵 $Y_e$，起始共有 $18$ 个实参数。令 Yukawa 耦合为零时，动能项具有
\begin{equation*}
 U(3)_{L_L}\times U(3)_{e_R}
\end{equation*}
的 $18$ 参数味基冗余。严格的最小标准模型中没有右手中微子，中微子保持无质量；把 $Y_e$ 对角化以后，三个轻子家族仍可分别做保持动能项和质量项不变的共同矢量式重相位。因此未破缺子群是 $U(1)^3$，维数为 $3$。可用于消去 $Y_e$ 坐标冗余的基变换参数数目为
\begin{equation*}
 18-3=15,
\end{equation*}
于是只剩
\begin{equation*}
 18-15=3
\end{equation*}
个带电轻子质量；在这个严格最小场内容中没有物理的 PMNS 混合角与相位。

所以全部 Yukawa 与味结构一共贡献
\begin{equation*}
 10+3=13
\end{equation*}
个物理参数。再加入三种规范耦合 $g_s,g,g'$、Higgs 势中的两个参数 $\mu_H^2,\lambda_H$，以及一个基不变的强相互作用 $CP$ 参数 $\bar\theta$，总数为
\begin{equation*}
 13+3+2+1=19.
\end{equation*}
于是得到\eqnrefs{eq:minimal-sm-parameter-count-dp11}。19 个连续参数就是耦合参数空间对场重定义轨道取商后的维数；粒子质量、混合参数和规范耦合只是这个商空间的一组物理坐标。
\end{derivation}

\begin{derivation}[从\eqnrefs{eq:sm-yukawa-general-dp11}到\eqnrefs{eq:higgs-fermion-coupling-dp11}]
对任意一种带电费米子 $f=u,d,e$，电弱破缺后正文 Yukawa 项都具有同一矩阵结构
\begin{equation*}
 \Lag_{Y,f}
 =-\frac{v_E}{\sqrt2}
 \left(\bar f_LY_ff_R+\bar f_RY_f^\dagger f_L\right)
 -\frac{\sqrt{\hbar c}\,h}{\sqrt2}
 \left(\bar f_LY_ff_R+\bar f_RY_f^\dagger f_L\right).
\end{equation*}
这里 $Y_f$ 是世代空间中的任意复 $N_g\times N_g$ 矩阵；它一般既不厄米也不正规，因此不能假定存在单个幺正矩阵 $U$ 使 $U^\dagger Y_fU$ 对角。

从两个厄米半正定矩阵
\begin{equation*}
 H_L\equiv Y_fY_f^\dagger,
 \qquad
 H_R\equiv Y_f^\dagger Y_f
\end{equation*}
出发。谱定理保证存在幺正矩阵 $U_{fL}$ 使
\begin{equation*}
 U_{fL}^\dagger H_LU_{fL}
 =\operatorname{diag}(y_1^2,\ldots,y_{N_g}^2),
 \qquad y_i\ge0.
\end{equation*}
设 $u_{Li}$ 是 $U_{fL}$ 的第 $i$ 个列向量，即
\begin{equation*}
 H_Lu_{Li}=y_i^2u_{Li}.
\end{equation*}
对任意 $y_i>0$，定义右奇异向量
\begin{equation*}
 u_{Ri}\equiv\frac1{y_i}Y_f^\dagger u_{Li}.
\end{equation*}
它们的内积为
\begin{align*}
 u_{Ri}^\dagger u_{Rj}
 &=\frac1{y_iy_j}
 u_{Li}^\dagger Y_fY_f^\dagger u_{Lj}\\
 &=\frac{y_j^2}{y_iy_j}u_{Li}^\dagger u_{Lj}
 =\delta_{ij},
\end{align*}
并且
\begin{align*}
 Y_fu_{Ri}
 &=\frac1{y_i}Y_fY_f^\dagger u_{Li}
 =y_i u_{Li},\\
 Y_f^\dagger u_{Li}
 &=y_i u_{Ri}.
\end{align*}
若存在 $y_i=0$，在 $\ker Y_f$ 与 $\ker Y_f^\dagger$ 中分别补齐正交基即可。把全部 $u_{Ri}$ 作为列向量组成幺正矩阵 $U_{fR}$，便有逐矩阵元关系
\begin{align*}
 (U_{fL}^\dagger Y_fU_{fR})_{ij}
 &=u_{Li}^\dagger Y_fu_{Rj}\\
 &=y_j u_{Li}^\dagger u_{Lj}
 =y_j\delta_{ij}.
\end{align*}
因此
\begin{equation*}
 \boxed{
 U_{fL}^\dagger Y_fU_{fR}
 =\operatorname{diag}(y_1,\ldots,y_{N_g})}
\end{equation*}
就是正文\eqnrefs{eq:yukawa-svd-dp11}。奇异值按定义非负，因此由 $m_f=v_E y_f/(\sqrt2 c)$ 得到的质量本征值可统一取非负实数。

定义新的左右手场
\begin{equation*}
 f_L'=U_{fL}^\dagger f_L,
 \qquad
 f_R'=U_{fR}^\dagger f_R.
\end{equation*}
因为 $U_{fL},U_{fR}$ 只作用在世代空间而且为常数矩阵，动能保持不变：
\begin{align*}
 \bar f_L\ii\hbar c\gamma^\mu D_\mu f_L
 &=\bar f_L'U_{fL}^\dagger
 \ii\hbar c\gamma^\mu D_\mu U_{fL}f_L'\\
 &=\bar f_L'\ii\hbar c\gamma^\mu D_\mu f_L',
\end{align*}
右手项同理。这里使用的是规范相互作用在同一电荷扇区的世代空间中正比于单位矩阵，因此 $D_\mu$ 与 $U_{fL,R}$ 对易。

质量项变为
\begin{align*}
 -\frac{v_E}{\sqrt2}\bar f_LY_ff_R+\mathrm{H.c.}
 &=-\frac{v_E}{\sqrt2}
 \bar f_L'U_{fL}^\dagger Y_fU_{fR}f_R'
 +\mathrm{H.c.}\\
 &=-\sum_i\frac{y_iv_E}{\sqrt2}
 \left(\bar f_{Li}'f_{Ri}'+\bar f_{Ri}'f_{Li}'\right)\\
 &=-\sum_i m_{fi}c^2\,\bar f_i'f_i',
\end{align*}
其中
\begin{equation*}
 \boxed{m_{fi}c^2=\frac{y_iv_E}{\sqrt2}},
\end{equation*}
即正文\eqnrefs{eq:fermion-masses-yukawa-dp11}。

关键点是 Higgs 涨落项含有\emph{完全相同的}矩阵 $Y_f$，因此不需要第二次对角化：
\begin{align*}
 \Lag_{hff}
 &=-\frac{\sqrt{\hbar c}\,h}{\sqrt2}
 \bar f_L'U_{fL}^\dagger Y_fU_{fR}f_R'
 +\mathrm{H.c.}\\
 &=-\sum_i
 \frac{y_i\sqrt{\hbar c}}{\sqrt2}
 h\left(\bar f_{Li}'f_{Ri}'+\bar f_{Ri}'f_{Li}'\right)\\
 &=-\sum_i
 \frac{m_{fi}c^2}{v_E}\sqrt{\hbar c}\,
 h\,\bar f_i'f_i'.
\end{align*}
这就是正文\eqnrefs{eq:higgs-fermion-coupling-dp11}。所以树级 Higgs 中性耦合在质量基中自动对角，且耦合强度与费米子质量线性成正比。

若某组奇异值严格简并，则对应简并子空间内仍允许额外幺正旋转而不改变质量矩阵；因此 CKM/PMNS 的独立参数计数需要单独注明“质量谱非简并”这一条件。质量本征值的定义和 Higgs 中性耦合的对角性本身并不依赖这一假设。
\end{derivation}

\begin{derivation}[反冲运动学与手征自旋求和]
记中微子散射角为 $\theta_\nu$，其末态能量为 $E'_\nu=E_\nu-T_e$。由 $P'^2=P^2=E_{e0}^2$、$P'=P+K-K'$ 得
\begin{align*}
 (K-K')^2&=-2E_{e0}T_e
 =-2E_\nu E'_\nu(1-\cos\theta_\nu),\\
 T_e&=\frac{E_\nu^2(1-\cos\theta_\nu)}
 {E_{e0}+E_\nu(1-\cos\theta_\nu)}.
\end{align*}
当 $\cos\theta_\nu=-1$ 时反冲最大，给\eqnrefs{eq:sm-nue-recoil-endpoint}。在这些运动学变量下，弱传播子展开的尺度比较为 $E_{e0}E_\nu/M_V^2\ll1$，而不是仅根据 $T_e/E_\nu$ 判断。

把\eqnrefs{eq:sm-nue-fermi-amplitude} 的两个电流分别平方，记
\begin{align*}
 L^{\mu\nu}&=\Tr[\slashed K'\gamma^\mu P_L\slashed K\gamma^\nu P_L],\\
 H_{LL,\mu\nu}&=\frac12\Tr[(\slashed P'+E_{e0})\gamma_\mu P_L
 (\slashed P+E_{e0})\gamma_\nu P_L],\\
 H_{RR,\mu\nu}&=\frac12\Tr[(\slashed P'+E_{e0})\gamma_\mu P_R
 (\slashed P+E_{e0})\gamma_\nu P_R].
\end{align*}
混合项 $H_{LR}+H_{RL}$ 由分别交换第二个投影得到。利用 $P_L\gamma^\mu=\gamma^\mu P_R$、$P_LP_R=0$ 和四个 Dirac 矩阵的迹，缩并结果为
\begin{align*}
 L^{\mu\nu}H_{LL,\mu\nu}&=8(P\cdot K)(P'\cdot K'),\\
 L^{\mu\nu}H_{RR,\mu\nu}&=8(P\cdot K')(P'\cdot K),\\
 L^{\mu\nu}(H_{LR}+H_{RL})_{\mu\nu}&=-8E_{e0}^2(K\cdot K').
\end{align*}
混合迹只留下两个质量插入；正是这一步产生负号，不能从两个独立手征概率相加得到。乘上顶角公共系数平方 $8[G_F^{(E)}]^2$ 后，
\begin{align*}
 \overline{|\mathcal M|^2}=64[G_F^{(E)}]^2\bigl[
 &g_L^2(P\cdot K)(P'\cdot K')+g_R^2(P\cdot K')(P'\cdot K)\\
 &-g_Lg_RE_{e0}^2(K\cdot K')\bigr].
\end{align*}
实验室不变量为
\begin{equation*}
 P\cdot K=P'\cdot K'=E_{e0}E_\nu,\quad
 P\cdot K'=P'\cdot K=E_{e0}(E_\nu-T_e),\quad
 K\cdot K'=E_{e0}T_e.
\end{equation*}
能量变量中的二体通量公式给
\begin{equation*}
 \frac{\dd\sigma}{\dd T_e}
 =\frac{(\hbar c)^2}{32\pi E_{e0}E_\nu^2}\overline{|\mathcal M|^2},
\end{equation*}
从而得到正文反冲谱。将中微子换成反中微子时，$L^{\mu\nu}$ 的反对称部分反号，两个同手征缩并互换而混合缩并不变，证明了 $g_L\leftrightarrow g_R$ 的替换规则。
\end{derivation}

\begin{derivation}[从微分截面到探测计数]
若电子靶数为 $N_e$、曝光时间为 $t_{\rm exp}$，入射味 $a$ 的微分通量为 $\phi_a(E_\nu)$，其单位为 $\mathrm{m^{-2}s^{-1}J^{-1}}$，则理想真实反冲分箱中的期望计数为
\begin{equation}
 N_{\rm bin}=N_et_{\rm exp}\sum_a\int\dd E_\nu\,\phi_a(E_\nu)
 \int_{T_1}^{\min(T_2,T_{\max})}\dd T_e\,
 \mathcal Q^{\rm det}(T_e)\frac{\dd\sigma_a}{\dd T_e}.
 \label{eq:sm-nue-event-folding}
\end{equation}
只对 $T_{\max}>T_1$ 的入射能量积分；其余贡献定义为零。逐项量纲为 $N_et_{\rm exp}\phi\dd E_\nu\sigma$，结果无量纲。若探测器报告重建能量而非真实动能，应再卷积归一化响应密度 $R(T_{\rm obs}\mid T_e)$，其单位为逆能量且满足 $\int\dd T_{\rm obs}R=1$。同一层级结构也用于本书后面的电子能谱反演：相互作用给真实谱，探测响应把真实谱映射成实际计数。
\end{derivation}

\begin{derivation}[从\eqnrefs{eq:eeff-chiral-coeff-dp11}到\eqnrefs{eq:eeff-spin-summed-dp11}]
考虑无质量过程 $e^-(p_1)e^+(p_2)\to f(k_1)\bar f(k_2)$，固定电子流手征 $i=L/R$ 与末态流手征 $j=L/R$。相应的流乘积为

\begin{equation*}
\mathcal J_{ij}
=[\bar v(p_2)\gamma^\mu P_i u(p_1)]
[\bar u(k_1)\gamma_\mu P_j v(k_2)].
\end{equation*}
对末态自旋求和并对初态自旋平均后，基本张量为
\begin{equation*}
\frac14\Sigma_{ij}
=\frac14\Tr(\slashed p_2\gamma^\mu P_i\slashed p_1\gamma^\nu P_i)
\Tr(\slashed k_1\gamma_\mu P_j\slashed k_2\gamma_\nu P_j).
\end{equation*}
利用
$P_L\slashed p=\slashed pP_R$、$P_R\slashed p=\slashed pP_L$ 以及
$P_R\gamma^\nu=\gamma^\nu P_L$，左手迹中的两个投影满足
\begin{align*}
 \Tr(\slashed a\gamma^\mu P_L\slashed b\gamma^\nu P_L)
 &=\Tr(\slashed a\gamma^\mu\slashed bP_R\gamma^\nu P_L),\\
 &=\Tr(\slashed a\gamma^\mu\slashed b\gamma^\nu P_L^2),\\
 &=\Tr(\slashed a\gamma^\mu\slashed b\gamma^\nu P_L).
\end{align*}
右手迹以 $L\leftrightarrow R$ 进行同样的投影代数。以左手投影为例，
\begin{align*}
\Tr(\slashed a\gamma^\mu\slashed b\gamma^\nu P_L)
&=\frac12\Tr(\slashed a\gamma^\mu\slashed b\gamma^\nu)
-\frac12\Tr(\slashed a\gamma^\mu\slashed b\gamma^\nu\gamma^5)\\
&=2\left[S^{\mu\nu}(a,b)+\ii A^{\mu\nu}(a,b)\right],
\end{align*}
其中
\begin{align*}
S^{\mu\nu}(a,b)
&\equiv a^\mu b^\nu+a^\nu b^\mu-\eta^{\mu\nu}a\cdot b,\\
A^{\mu\nu}(a,b)
&\equiv\epsilon^{\alpha\mu\beta\nu}a_\alpha b_\beta.
\end{align*}
右手投影只把 $A^{\mu\nu}$ 的符号反转。现在把“手征选择”真正做成张量收缩。由于一个张量对称、另一个反对称，
\begin{equation*}
S^{\mu\nu}(a,b)A_{\mu\nu}(c,d)=0.
\end{equation*}
对称部分直接展开为
\begin{equation*}
S^{\mu\nu}(a,b)S_{\mu\nu}(c,d)
=2\left[(a\cdot c)(b\cdot d)+(a\cdot d)(b\cdot c)\right].
\end{equation*}
反对称部分使用
\begin{equation*}
\epsilon^{\alpha\mu\beta\nu}\epsilon_{\gamma\mu\delta\nu}
=-2\left(\delta^\alpha_{\gamma}\delta^\beta_{\delta}
-\delta^\alpha_{\delta}\delta^\beta_{\gamma}\right)
\end{equation*}
得到
\begin{equation*}
(\ii A^{\mu\nu})(\ii A_{\mu\nu})
=2\left[(a\cdot d)(b\cdot c)-(a\cdot c)(b\cdot d)\right].
\end{equation*}
因此同手征时对称部分与反对称部分相加，只保留第二种配对；异手征时其中一个 $A$ 反号，只保留第一种配对。取
$(a,b,c,d)=(p_2,p_1,k_1,k_2)$，并记住每个手征迹前的因子 $2$ 以及初态平均 $1/4$，便有
\begin{align*}
\frac14\Sigma_{LL}
=\frac14\Sigma_{RR}
&=4(p_1\cdot k_1)(p_2\cdot k_2),\\
\frac14\Sigma_{LR}
=\frac14\Sigma_{RL}
&=4(p_1\cdot k_2)(p_2\cdot k_1).
\end{align*}

在质心系令 $x=\cos\theta$，无质量运动学给出
\begin{align*}
2p_1\cdot k_1&=-t=\frac{s}{2}(1-x),\\
2p_1\cdot k_2&=-u=\frac{s}{2}(1+x),
\end{align*}
并且由质心系的交换对称性
\begin{equation*}
p_2\cdot k_2=p_1\cdot k_1,
\qquad
p_2\cdot k_1=p_1\cdot k_2.
\end{equation*}
于是两类手征块分别产生 $(1-x)^2$ 与 $(1+x)^2$。按照 $C_{LL},C_{LR},C_{RL},C_{RR}$ 的外流手征约定重新排列后，即得到\eqnrefs{eq:eeff-spin-summed-dp11}。因此二次角权重由手征投影下的对称四伽马迹与反对称 Levi--Civita 迹共同确定。
\end{derivation}

\begin{derivation}[从\eqnrefs{eq:pmns-definition-dp11}到\eqnrefs{eq:neutrino-oscillation-probability-dp14}]
Takagi 分解后，味本征场与质量本征场的左手分量满足

\begin{equation*}
\nu_{\alpha L}=\sum_{i=1}^3U_{\alpha i}\nu_{iL},
\qquad U\equiv U_{\rm PMNS}.
\end{equation*}
因此在产生点制备的味态为
\begin{equation*}
|\nu_\alpha\rangle=\sum_iU_{\alpha i}^*|\nu_i\rangle.
\end{equation*}

各质量分量具有共同主导能量 $E$ 且满足 $m_i^2c^4/E^2\ll1$ 时，精确色散关系
\begin{equation*}
E_i^2=p_i^2c^2+m_i^2c^4
\end{equation*}
给出
\begin{align*}
p_ic
&=E\sqrt{1-\frac{m_i^2c^4}{E^2}}\\
&=E\left[1-\frac{m_i^2c^4}{2E^2}
+\mathcal O\!\left(\frac{m_i^4c^8}{E^4}\right)\right],
\end{align*}
所以
\begin{equation*}
p_i=\frac Ec-\frac{m_i^2c^3}{2E}
+\mathcal O\!\left(\frac{m_i^4c^7}{E^3}\right).
\end{equation*}
传播距离 $L$ 的平面波相位为
\begin{equation*}
\exp\!\left[-\frac{\ii}{\hbar}(E_it-p_iL)\right].
\end{equation*}
在各波包仍相干重叠时取 $t=L/c$ 到主导阶，并去掉所有质量分量共有的整体相位，留下
\begin{equation*}
\exp\!\left[-\ii\frac{m_i^2c^3L}{2E\hbar}\right].
\end{equation*}
于是
\begin{equation*}
\mathcal A_{\alpha\to\beta}(L)
=\sum_iU_{\beta i}U_{\alpha i}^*
\exp\!\left[-\ii\frac{m_i^2c^3L}{2E\hbar}\right].
\end{equation*}

取模平方并把 $i=j$ 与 $i\neq j$ 分开：
\begin{align*}
P_{\alpha\to\beta}
={}&\sum_i|U_{\beta i}|^2|U_{\alpha i}|^2\\
&+2\Re\sum_{i<j}
U_{\beta i}U_{\alpha i}^*U_{\beta j}^*U_{\alpha j}
\exp\!\left[-\ii\frac{\Delta m_{ij}^2c^3L}{2E\hbar}\right].
\end{align*}
利用幺正性 $\sum_iU_{\beta i}U_{\alpha i}^*=\delta_{\alpha\beta}$，并使用 $K_{ij}^{\alpha\beta}$ 与 $\Delta_{ij}$ 的定义，将指数的实部和虚部分开，即得到\eqnrefs{eq:neutrino-oscillation-probability-dp14}。

\end{derivation}

\begin{derivation}[从\eqnrefs{eq:smeft-geometric-normalization-dp13}到\eqnrefs{eq:heavy-scalar-ch-loop-dp32}]
取几何量纲为 $\mathrm{m^{-1}}$ 的重实标量 $S$，并令

\begin{equation*}
\widehat\Lag_S
=\frac12(\partial_\mu S)(\partial^\mu S)
-\frac12\Lambda_*^2S^2
-\frac\kappa2S^2(\widehat\Phi^\dagger\widehat\Phi),
\end{equation*}
其中 $\kappa$ 无量纲。经典分支 $S=0$ 不产生 $\mathcal O_H$，所以最低非零项来自高斯泛函积分的一圈行列式。转到欧几里得形式并令 $\widehat H=\widehat\Phi^\dagger\widehat\Phi$：
\begin{equation*}
\frac{\Delta\Gamma_E^{(1)}}{\hbar}
=\frac12\Tr\ln[-\partial_E^2+\Lambda_*^2+\kappa\widehat H].
\end{equation*}
在缓慢变化背景下忽略 $\partial\widehat H$，把场依赖部分写成
\begin{equation*}
\frac12\int\frac{\dd^4 k_E}{(2\pi)^4}
\ln\!\left(1+\frac{\kappa\widehat H}{k_E^2+\Lambda_*^2}\right).
\end{equation*}
展开 $\ln(1+u)=u-u^2/2+u^3/3+\cdots$。前两项分别重整化 Higgs 质量和四次耦合，首个新的六维结构来自 $u^3$：
\begin{equation*}
\Delta\widehat V_{E,H^3}^{(1)}
=\frac{\kappa^3\widehat H^3}{6}I_3(\Lambda_*),
\end{equation*}
其中
\begin{equation*}
I_3(\Lambda_*)
=\int\frac{\dd^4 k_E}{(2\pi)^4}\frac1{(k_E^2+\Lambda_*^2)^3}.
\end{equation*}
用 Schwinger 参数表示
\begin{equation*}
\frac1{A^3}=\frac12\int_0^\infty\dd\tau\,\tau^2\ee^{-\tau A}
\end{equation*}
以及四维高斯积分
\begin{equation*}
\int\frac{\dd^4 k_E}{(2\pi)^4}\ee^{-\tau k_E^2}
=\frac1{16\pi^2\tau^2},
\end{equation*}
得到
\begin{align*}
I_3(\Lambda_*)
&=\frac1{32\pi^2}\int_0^\infty\dd\tau\,\ee^{-\tau\Lambda_*^2}\\
&=\frac1{32\pi^2\Lambda_*^2}.
\end{align*}
闵可夫斯基拉格朗日密度中的势能项带负号，因此
\begin{equation*}
\Delta\widehat\Lag_{6,H}^{(1)}
=-\frac{\kappa^3}{192\pi^2\Lambda_*^2}
(\widehat\Phi^\dagger\widehat\Phi)^3.
\end{equation*}
与\eqnrefs{eq:smeft-geometric-normalization-dp13} 比较，即得到\eqnrefs{eq:heavy-scalar-ch-loop-dp32}。改回 SI Higgs 场后，系数前的整体量纲仍是 $\mathrm{J\,m^{-3}}$。
\end{derivation}



\begin{derivation}[从\eqnrefs{eq:weinberg-bare-renorm-dp91}到\eqnrefs{eq:weinberg-matrix-rg-dp32}]
计算两个左手轻子和两个 Higgs 的 1PI 四点函数。取 $d=4-2\epsilon$，并保留一个非零外部不变量来分开紫外与红外奇点。一圈两传播子子图都含公共泡积分
\begin{equation*}
 B(p^2)=\mu_{\rm DR}^{2\epsilon}
 \int\frac{\dd^d\ell}{(2\pi)^d}
 \frac{1}{(\ell^2+\ii0)[(\ell+p)^2+\ii0]}.
\end{equation*}
用 $1/(AB)=\int_0^1\dd x\,[xA+(1-x)B]^{-2}$，再令 $k=\ell+(1-x)p$，有
\begin{align*}
 B(p^2)
 &=\mu_{\rm DR}^{2\epsilon}\int_0^1\dd x
 \int\frac{\dd^dk}{(2\pi)^d}
 \frac{1}{[k^2-x(1-x)p^2+\ii0]^2}\\
 &=\frac{\ii}{(4\pi)^{2-\epsilon}}\Gamma(\epsilon)
 \int_0^1\dd x
 \left[\frac{\mu_{\rm DR}^2}{-x(1-x)p^2-\ii0}\right]^\epsilon.
\end{align*}
因此
\begin{equation*}
 B(p^2)\big|_{\rm UV}
 =\frac{\ii}{16\pi^2}\Delta_\epsilon,
 \qquad
 \Delta_\epsilon\equiv\frac1\epsilon-\gamma_E+\ln4\pi.
\end{equation*}
所有一圈反项都由这个极点乘上不同的味、弱同位旋与规范群系数。

味矩阵的左右位置必须由显式指标固定。带电轻子 Yukawa 项是 $\bar L_{Li}(Y_e)_{ia}\Phi e_{Ra}$，所以 Weinberg 算符左手轻子腿上的自能味矩阵是
\begin{align*}
 (\mathcal Y_L)_{ij}
 &=\sum_a(Y_e)_{ia}(Y_e^\dagger)_{aj}\\
 &=(Y_eY_e^\dagger)_{ij}.
\end{align*}
若 Yukawa 子图接在系数矩阵右侧轻子腿，味链为
\begin{align*}
 \sum_{a,\rho}C^{(5)}_{i\rho}(Y_e)_{\rho a}(Y_e^\dagger)_{aj}
 &=\sum_\rho C^{(5)}_{i\rho}(\mathcal Y_L)_{\rho j}\\
 &=(C^{(5)}\mathcal Y_L)_{ij};
\end{align*}
接在电荷共轭的左侧轻子腿时则为
\begin{align*}
 \sum_{a,\rho}(Y_e^*)_{ia}(Y_e^T)_{a\rho}C^{(5)}_{\rho j}
 &=\sum_\rho(\mathcal Y_L^T)_{i\rho}C^{(5)}_{\rho j}\\
 &=(\mathcal Y_L^TC^{(5)})_{ij}.
\end{align*}
这一步固定了 RGE 中矩阵乘法的左右次序；不能把 $Y_eY_e^\dagger$ 换成作用在右手单态空间的 $Y_e^\dagger Y_e$。

用公共因子 $(16\pi^2\epsilon)^{-1}$ 表示一圈简单极点。四点 1PI 顶角的局域反项逐类相加为
\begin{equation*}
 \delta V_5^{(1)}
 =\frac{1}{16\pi^2\epsilon}
 \left[
 \left(2\lambda_H+\frac14g'^2-\frac34g^2\right)C^{(5)}
 -\mathcal Y_L^TC^{(5)}-C^{(5)}\mathcal Y_L
 \right].
\end{equation*}
其中最后两项正是上面逐指标得到的两条 Yukawa 味链。$2\lambda_H$ 来自四 Higgs 顶角与 Weinberg 顶角的两个 Higgs 指标收缩；$g'^2/4$ 与 $-3g^2/4$ 分别来自 $U(1)_Y$ 与 $SU(2)_L$ 的规范交换及顶角子图。

还必须把四条外腿的场强反项并入同一个裸系数。左手轻子二重态和 Higgs 二重态的一圈波函数反项为
\begin{align*}
 \delta Z_L^{(1)}
 &=-\frac{1}{16\pi^2\epsilon}\frac14
 \left[g'^2+3g^2+2\mathcal Y_L\right],\\
 \delta Z_\Phi^{(1)}
 &=\frac{1}{16\pi^2\epsilon}\frac12
 \left[g'^2+3g^2-2\mathcal T_Y\right].
\end{align*}
第一式的矩阵部分仍作用在左手轻子味空间；第二式中的闭合费米子环已把味与颜色取迹，因此只剩标量 $\mathcal T_Y$。

把
\begin{equation*}
 Z_L^{-1/2}=I-\frac12\delta Z_L^{(1)}+O(g^4,Y^4),
 \qquad
 Z_\Phi^{-1}=1-\delta Z_\Phi^{(1)}+O(g^4,Y^4)
\end{equation*}
代入正文\eqnrefs{eq:weinberg-bare-renorm-dp91}。一圈 Wilson 系数反项为
\begin{equation*}
 \delta C_5^{(1)}
 =\delta V_5^{(1)}
 -\frac12(\delta Z_L^{(1)})^T C^{(5)}
 -\frac12C^{(5)}\delta Z_L^{(1)}
 -\delta Z_\Phi^{(1)}C^{(5)}.
\end{equation*}
把顶角与外腿反项按群结构逐项相加。$U(1)_Y$ 部分为
\begin{equation*}
 \frac14g'^2
 +\frac18g'^2+\frac18g'^2
 -\frac12g'^2=0,
\end{equation*}
所以一圈 Weinberg RGE 中没有 $g'^2C^{(5)}$ 项。$SU(2)_L$ 部分为
\begin{equation*}
 -\frac34g^2
 +\frac38g^2+\frac38g^2
 -\frac32g^2
 =-\frac32g^2.
\end{equation*}
带电轻子非普适部分则为
\begin{align*}
 &-\left(\mathcal Y_L^TC^{(5)}+C^{(5)}\mathcal Y_L\right)
 +\frac14\left(\mathcal Y_L^TC^{(5)}+C^{(5)}\mathcal Y_L\right)\\
 &\hspace{28mm}
 =-\frac34\left(\mathcal Y_L^TC^{(5)}+C^{(5)}\mathcal Y_L\right),
\end{align*}
而 Higgs 波函数中的 $-2\mathcal T_Y$ 经 $-\delta Z_\Phi$ 给出 $+\mathcal T_YC^{(5)}$。因此
\begin{align*}
 \delta C_5^{(1)}
 =\frac{1}{16\pi^2\epsilon}
 \Bigg\{&
 \left(2\lambda_H-\frac32g^2+\mathcal T_Y\right)C^{(5)}\\
 &-\frac34\left(\mathcal Y_L^TC^{(5)}+C^{(5)}\mathcal Y_L\right)
 \Bigg\}.
\end{align*}

简单极点转换为 beta 函数时，$A_1$ 不能当作与尺度无关的常数，因为在 $d=4-2\epsilon$ 中各无量纲耦合具有不同的工程维数。对任意一组重整化参数 $\kappa_i$，把裸量统一写成
\begin{equation*}
 \kappa_{i,B}=\mu_{\rm DR}^{\rho_i\epsilon}
 \left[\kappa_i+\frac{a_i^{(1)}(\boldsymbol\kappa)}{\epsilon}
 +\sum_{n\ge2}\frac{a_i^{(n)}(\boldsymbol\kappa)}{\epsilon^n}\right].
\end{equation*}
固定所有裸量并对 $t=\ln\mu_{\rm DR}$ 求导。$d$ 维 beta 函数分成工程维数与四维部分，
\begin{equation*}
 \frac{\dd\kappa_i}{\dd t}
 =-\rho_i\epsilon\kappa_i+\beta_i(\boldsymbol\kappa).
\end{equation*}
代回裸量的导数，在一圈阶只需保留 $a_i^{(1)}$ 中各耦合的工程维数项：
\begin{align*}
0={}&\rho_i\epsilon
 \left(\kappa_i+\frac{a_i^{(1)}}{\epsilon}\right)
 -\rho_i\epsilon\kappa_i+\beta_i
 +\frac1\epsilon\sum_j
 \left(-\rho_j\epsilon\kappa_j\right)
 \frac{\partial a_i^{(1)}}{\partial\kappa_j}
 +O((16\pi^2)^{-2})\\
={}&\beta_i+\rho_i a_i^{(1)}
 -\sum_j\rho_j\kappa_j
 \frac{\partial a_i^{(1)}}{\partial\kappa_j}
 +O((16\pi^2)^{-2}).
\end{align*}
因此 MS 类方案的一圈 simple-pole 公式是
\begin{equation*}
 \beta_i^{(1)}=
 \sum_j\rho_j\kappa_j
 \frac{\partial a_i^{(1)}}{\partial\kappa_j}
 -\rho_i a_i^{(1)}.
\end{equation*}
对 Weinberg 系数，$\rho_{C_5}=2$；当前 Higgs 势归一化下
$\rho_{\lambda_H}=2$、$\rho_g=\rho_{g'}=1$、$\rho_Y=1$。把上面的简单极点残数写成
\begin{equation*}
 a_{C_5}^{(1)}=\frac1{16\pi^2}\left[
 \left(2\lambda_H-\frac32g^2+\mathcal T_Y\right)C^{(5)}
 -\frac34\left(\mathcal Y_L^TC^{(5)}+C^{(5)}\mathcal Y_L\right)
 \right].
\end{equation*}
它对 $C^{(5)}$ 是一次齐次的。把矩阵元 $C^{(5)}_{ij}$ 视为独立坐标，Euler 定理给出
\begin{equation*}
 2\sum_{ij}C^{(5)}_{ij}
 \frac{\partial a_{C_5}^{(1)}}{\partial C^{(5)}_{ij}}
 -2a_{C_5}^{(1)}=0.
\end{equation*}
Wilson 系数自身的工程维数项因而正好抵消，剩下只需对耦合常数作齐次计数。四 Higgs 耦合项满足
\begin{equation*}
 2\lambda_H\frac{\partial}{\partial\lambda_H}
 \left(2\lambda_H C^{(5)}\right)=4\lambda_H C^{(5)},
\end{equation*}
规范项满足
\begin{equation*}
 g\frac{\partial}{\partial g}
 \left(-\frac32g^2C^{(5)}\right)=-3g^2C^{(5)}.
\end{equation*}
$\mathcal T_Y$ 与 $\mathcal Y_L$ 都是 Yukawa 矩阵的二次齐次函数；把 $Y_f$ 与 $Y_f^*$ 视为独立复变量作 Euler 计数，得到
\begin{align*}
 \sum_f\left(Y_f\!:\frac{\partial}{\partial Y_f}
 +Y_f^*\!:\frac{\partial}{\partial Y_f^*}\right)\mathcal T_Y
 &=2\mathcal T_Y,\\
 \left(Y_e\!:\frac{\partial}{\partial Y_e}
 +Y_e^*\!:\frac{\partial}{\partial Y_e^*}\right)\mathcal Y_L
 &=2\mathcal Y_L.
\end{align*}
于是 simple-pole 公式逐项给出
\begin{align*}
 16\pi^2\mu_{\rm DR}\frac{\dd C^{(5)}}{\dd\mu_{\rm DR}}
 ={}&\left(4\lambda_H-3g^2+2\mathcal T_Y\right)C^{(5)}\\
 &-\frac32\left(\mathcal Y_L^TC^{(5)}+C^{(5)}\mathcal Y_L\right),
\end{align*}
即正文\eqnrefs{eq:weinberg-matrix-rg-dp32}。右端在 $C^{(5)T}=C^{(5)}$ 时仍为复对称矩阵，所以 Majorana 系数的对称性在一圈运行下保持不变。
\end{derivation}

\begin{derivation}[从\eqnrefs{eq:delta-r-definition-dp11}到\eqnrefs{eq:sm-mw-rho-response-dp95}]
把\eqnrefs{eq:delta-r-definition-dp11} 改写为
\begin{equation*}
M_W^2s_W^2=A(1+\Delta r),
\qquad
A\equiv\frac{\pi\alpha}{\sqrt2G_F^{(E)}}.
\end{equation*}
在壳方案中 $s_W^2=1-M_W^2/M_Z^2$，因此定义
\begin{equation*}
f(x)\equiv x\left(1-\frac{x}{M_Z^2}\right),
\qquad x\equiv M_W^2,
\end{equation*}
则精密输入关系就是 $f(x)=A(1+\Delta r)$。树级点 $x_0$ 满足 $f(x_0)=A$。把精确解写成 $x=x(\Delta r)$，满足
\begin{equation*}
 f[x(\Delta r)]=A(1+\Delta r),
 \qquad x(0)=x_0.
\end{equation*}
对 $\Delta r$ 在零点求导，隐函数定理直接给
\begin{equation*}
 f'(x_0)\left.\frac{\dd x}{\dd(\Delta r)}\right|_0=A.
\end{equation*}
而
\begin{equation*}
f'(x_0)
=1-\frac{2M_{W,0}^2}{M_Z^2}
=s_{W,0}^2-c_{W,0}^2,
\qquad
A=x_0s_{W,0}^2.
\end{equation*}
因此平方质量对 $\Delta r$ 的精确线性响应系数为
\begin{equation*}
\left.\frac1{x}\frac{\dd x}{\dd(\Delta r)}\right|_0
=-\frac{s_{W,0}^2}{c_{W,0}^2-s_{W,0}^2}.
\end{equation*}
又因为 $M_W=\sqrt{x}$，链式法则给
\begin{equation*}
\boxed{
\left.\frac1{M_W}\frac{\dd M_W}{\dd(\Delta r)}\right|_0
=-\frac{s_{W,0}^2}{2(c_{W,0}^2-s_{W,0}^2)}.}
\end{equation*}
这就是\eqnrefs{eq:sm-mw-linear-response-dp95} 所表示的一圈线性响应系数；有限 $\Delta r$ 的非线性位移由精确二次方程 $f(x)=A(1+\Delta r)$ 决定。

接着把 $\Delta\rho$ 在 $\Delta r$ 中的系数从在壳反项显式抽出来。树级带电流振幅含组合 $g^2/M_W^2=g_{\rm em}^2/(s_W^2M_W^2)$，因此参数反项中由弱混合角产生
\begin{equation*}
\frac{\delta(g^2/M_W^2)}{g^2/M_W^2}
\supset -2\frac{\delta s_W}{s_W}.
\end{equation*}
而正文\eqnrefs{eq:sw-counterterm-massrelation-dp11} 给
\begin{equation*}
-2\frac{\delta s_W}{s_W}
=\frac{c_W^2}{s_W^2}
\left(\frac{\delta m_W^2}{m_W^2}-\frac{\delta m_Z^2}{m_Z^2}\right).
\end{equation*}
在零动量的守护破缺部分，质量反项差与重整化自能差具有相反号；按正文\eqnrefs{eq:delta-rho-definition-dp11} 的定义，
\begin{equation*}
\left(\frac{\delta m_W^2}{m_W^2}-\frac{\delta m_Z^2}{m_Z^2}\right)_{\rm custodial}
=-\Delta\rho.
\end{equation*}
所以 $\Delta r$ 中这部分必为
\begin{equation*}
\left.\Delta r\right|_{\Delta\rho}
=-\frac{c_W^2}{s_W^2}\Delta\rho.
\end{equation*}
再把光子真空极化主导的电荷运行记为 $\Delta\alpha$，其余有限自能、顶角、盒图与反项归入 $\Delta r_{\rm rem}$，得到
\begin{equation*}
\Delta r
=\Delta\alpha-\frac{c_W^2}{s_W^2}\Delta\rho+\Delta r_{\rm rem}.
\end{equation*}
因此 $-c_W^2/s_W^2$ 不是经验分解系数，而是由在壳定义 $s_W^2=1-m_W^2/m_Z^2$ 对带电流振幅的参数依赖直接固定。
把 $\Delta r$ 中与 $\Delta\rho$ 成线性的部分代入上面的响应系数，得到
\begin{align*}
\left.\frac{\delta M_W}{M_W}\right|_{\Delta\rho}
&=-\frac{s_W^2}{2(c_W^2-s_W^2)}
\left(-\frac{c_W^2}{s_W^2}\Delta\rho\right)\\
&=\frac{c_W^2}{2(c_W^2-s_W^2)}\Delta\rho,
\end{align*}
得到\eqnrefs{eq:sm-mw-rho-response-dp95}。再代入\eqnrefs{eq:delta-rho-top-dp11} 就得到正文的重顶响应\eqnrefs{eq:sm-mw-top-response-dp95}。
\end{derivation}

\begin{derivation}[从\eqnrefs{eq:sm-one-loop-determinant-master-dp68}到\eqnrefs{eq:sm-coleman-weinberg-si-dp12}]
正文\eqnrefs{eq:sm-one-loop-determinant-master-dp68} 已把常数 Higgs 背景上的二次涨落统一为函数行列式。先取一个实玻色自由度，记其背景质量能量平方为 $\mathcal M^2$：
\begin{equation*}
 V_1^{(B)}(\mathcal M^2)
 =\frac{\hbar c}{2}
 \left(\frac{\mu_{\rm DR}}{\hbar}\right)^{2\epsilon}
 \int\frac{\dd^dk}{(2\pi)^d}
 \ln\!\left[k^2+\frac{\mathcal M^2}{(\hbar c)^2}\right],
 \qquad d=4-2\epsilon.
\end{equation*}
直接积分对数不方便；先对 $\mathcal M^2$ 求导：
\begin{equation*}
 \frac{\partial V_1^{(B)}}{\partial\mathcal M^2}
 =\frac{1}{2\hbar c}
 \left(\frac{\mu_{\rm DR}}{\hbar}\right)^{2\epsilon}
 \int\frac{\dd^dk}{(2\pi)^d}
 \frac{1}{k^2+\mathcal M^2/(\hbar c)^2}.
\end{equation*}
欧氏母积分
\begin{equation*}
 \int\frac{\dd^dk}{(2\pi)^d}\frac1{k^2+m^2}
 =\frac{\Gamma(1-d/2)}{(4\pi)^{d/2}}(m^2)^{d/2-1}
\end{equation*}
给出
\begin{equation*}
 \frac{\partial V_1^{(B)}}{\partial\mathcal M^2}
 =\frac{\mathcal M^2}{2(4\pi)^2(\hbar c)^3}
 \Gamma(-1+\epsilon)(4\pi)^\epsilon
 \left(\frac{E_R^2}{\mathcal M^2}\right)^\epsilon,
 \qquad E_R=c\mu_{\rm DR}.
\end{equation*}
把 $\mathcal M^2$ 从零积分到当前背景值；与背景无关的积分常数只平移真空能，可吸收到宇宙学常数型反项。利用
\begin{align*}
 \Gamma(-2+\epsilon)
 &=\frac1{2\epsilon}+\frac34-\frac{\gamma_E}{2}+\order(\epsilon),\\
 (4\pi)^\epsilon
 &=1+\epsilon\ln4\pi+\order(\epsilon^2),\\
 \left(\frac{E_R^2}{\mathcal M^2}\right)^\epsilon
 &=1+\epsilon\ln\frac{E_R^2}{\mathcal M^2}+\order(\epsilon^2),
\end{align*}
可把未减除结果整理为
\begin{equation*}
 V_1^{(B)}
 =\frac{\mathcal M^4}{64\pi^2(\hbar c)^3}
 \left[
 -\frac1{\bar\epsilon}
 +\ln\frac{\mathcal M^2}{E_R^2}-\frac32
 \right]+\order(\epsilon),
 \qquad
 \frac1{\bar\epsilon}\equiv\frac1\epsilon-\gamma_E+\ln4\pi.
\end{equation*}
$\overline{\rm MS}$ 反项只减去 $1/\bar\epsilon$，所以每个实标量自由度留下
\begin{equation*}
 V_{1,\rm ren}^{(B)}
 =\frac{\mathcal M^4}{64\pi^2(\hbar c)^3}
 \left(\ln\frac{\mathcal M^2}{E_R^2}-\frac32\right).
\end{equation*}

费米高斯积分产生 $\det\mathcal D_F$ 而不是 $\det^{-1/2}\mathcal D_B$，因此闭合费米子环带负号；Dirac 自旋、颜色等多重度再乘入 $n_i$。质量矢量场的常数 $5/6$ 还不能简单从标量结果“照抄”。在维数正则化中，横向 massive-vector 行列式的极化数是 $d-1=3-2\epsilon$，因此其裸贡献含
\begin{align*}
 (3-2\epsilon)
 \left[-\frac1{\bar\epsilon}
 +\ln\frac{\mathcal M_V^2}{E_R^2}-\frac32\right]
 &=-\frac3{\bar\epsilon}
 +3\ln\frac{\mathcal M_V^2}{E_R^2}
 -\frac92+2+\order(\epsilon)\\
 &=-\frac3{\bar\epsilon}
 +3\left[
 \ln\frac{\mathcal M_V^2}{E_R^2}-\frac56
 \right]+\order(\epsilon).
\end{align*}
这里有限的 $+2$ 正是 $(-2\epsilon)$ 乘紫外极点产生的项；它把每个物理矢量自由度的常数从 $3/2$ 改为 $5/6$。在正文采用的 Landau 规范组织下，Goldstone 与 ghost 的背景依赖另外计入各自的 $n_i,\mathcal M_i$，不会再重复计入这一横向矢量多重度。

把所有粒子种类统一写成带符号的 $n_i$，并加回树级势，最终得到
\begin{equation*}
 V_{\rm eff}=V_0+\frac1{64\pi^2(\hbar c)^3}
 \sum_i n_i\mathcal M_i^4
 \left[\ln\frac{\mathcal M_i^2}{E_R^2}-c_i\right]
 +\order(\hbar^2),
\end{equation*}
其中标量和费米子 $c_i=3/2$，横向质量矢量 $c_i=5/6$，即正文\eqnrefs{eq:sm-coleman-weinberg-si-dp12}。每项的 $\mathcal M_i^4/(\hbar c)^3$ 均具有 $\mathrm{J\,m^{-3}}$ 量纲。
\end{derivation}

\begin{derivation}[从\eqnrefs{eq:sm-higgs-vector-rules-dp43}到\eqnrefs{eq:higgs-to-vv-rates-dp11}]
正文\eqnrefs{eq:sm-higgs-vector-rules-dp43} 给出 $hWW$ 与 $hZZ$ 顶角。对两个有质量矢量玻色子的三个物理极化求和，利用在壳二体运动学整理纵向贡献，再乘二体相空间并对 $ZZ$ 加入相同粒子因子，即得到正文\eqnrefs{eq:higgs-to-vv-rates-dp11}。

记 $V=W$ 或 $Z$，并定义
\begin{equation*}
M_V\equiv m_Vc^2,
\qquad
x_V\equiv\frac{m_V^2}{m_h^2}.
\end{equation*}
质量本征基的 Higgs--矢量顶角具有
\begin{equation*}
\ii g_{hVV}\eta_{\mu\nu},
\qquad
g_{hVV}=\frac{2M_V^2}{v_E}.
\end{equation*}
因此固定两条末态极化时
\begin{equation*}
\mathcal M_{\lambda_1\lambda_2}
=g_{hVV}\,
\epsilon_{1}^{(\lambda_1)*}\!\cdot
\epsilon_{2}^{(\lambda_2)*}.
\end{equation*}
每个在壳有质量矢量具有三个物理极化，完备关系为
\begin{equation*}
P_{\mu\nu}(k)
\equiv
\sum_\lambda
\epsilon_\mu^{(\lambda)}(k)
\epsilon_\nu^{(\lambda)*}(k)
=-\eta_{\mu\nu}
+\frac{k_\mu k_\nu}{m_V^2c^2}.
\end{equation*}
于是极化求和后的矩阵元平方为
\begin{align*}
\sum_{\lambda_1,\lambda_2}|\mathcal M|^2
&=g_{hVV}^2
\eta^{\mu\nu}\eta^{\rho\sigma}
P^{(1)}_{\mu\rho}P^{(2)}_{\nu\sigma}\\
&=g_{hVV}^2
\left[
4-\frac{k_1^2}{m_V^2c^2}
-\frac{k_2^2}{m_V^2c^2}
+\frac{(k_1\!\cdot k_2)^2}{m_V^4c^4}
\right].
\end{align*}
两条末态均在壳，故 $k_1^2=k_2^2=m_V^2c^2$，从而
\begin{equation*}
\sum|\mathcal M|^2
=g_{hVV}^2
\left[2+\frac{(k_1\!\cdot k_2)^2}{m_V^4c^4}\right].
\end{equation*}
Higgs 质量壳满足 $(k_1+k_2)^2=m_h^2c^2$，因此
\begin{align*}
2k_1\!\cdot k_2
&=m_h^2c^2-k_1^2-k_2^2\\
&=(m_h^2-2m_V^2)c^2,
\end{align*}
即
\begin{equation*}
\frac{(k_1\!\cdot k_2)^2}{m_V^4c^4}
=\frac{(1-2x_V)^2}{4x_V^2}.
\end{equation*}
代回后
\begin{align*}
2+\frac{(1-2x_V)^2}{4x_V^2}
&=\frac{8x_V^2+1-4x_V+4x_V^2}{4x_V^2}\\
&=\frac{1-4x_V+12x_V^2}{4x_V^2}.
\end{align*}
再乘 $g_{hVV}^2=4M_V^4/v_E^2$，得到熟悉的运动学多项式
\begin{equation*}
\sum|\mathcal M|^2
=\frac{(m_hc^2)^4}{v_E^2}
\left(1-4x_V+12x_V^2\right).
\end{equation*}
二体相空间的末态速度因子为
\begin{equation*}
\beta_V=\sqrt{1-4x_V}.
\end{equation*}
把该矩阵元平方、$\beta_V$ 和 $1/v_E^2=\sqrt2G_F^{(E)}$ 代入一体衰变母式，得到\eqnrefs{eq:higgs-to-vv-rates-dp11}。$ZZ$ 末态的两粒子完全相同，因此其相空间要除以 $2!$；$W^+W^-$ 可区分，不含这一因子，所以两个宽度的整体前因子相差二倍。
\end{derivation}

\begin{derivation}[从\eqnrefs{eq:ew-onshell-mass-field-ct-dp11}到\eqnrefs{eq:sw-counterterm-massrelation-dp11}]
横向重整化逆传播子写成
\begin{equation*}
 [D_{V,T}(p^2)]^{-1}
 =p^2-m_V^2c^2+\Sigma_{VV,T}^{\rm loop}(p^2)
 -\delta m_V^2c^2+(p^2-m_V^2c^2)\delta Z_V.
\end{equation*}
极点仍位于 $p^2=m_V^2c^2$ 且留数为一，等价于
\begin{align*}
 0&=\Re\widehat\Sigma_{VV,T}(m_V^2c^2)
 \quad\Longrightarrow\quad
 \delta m_V^2c^2=\Re\Sigma_{VV,T}^{\rm loop}(m_V^2c^2),\\
 0&=\Re\widehat\Sigma'_{VV,T}(m_V^2c^2)
 \quad\Longrightarrow\quad
 \delta Z_V=-\Re\Sigma_{VV,T}^{{\rm loop}\,'}(m_V^2c^2),
\end{align*}
即正文\eqnrefs{eq:ew-onshell-mass-field-ct-dp11}。

在壳定义 $s_W^2=1-m_W^2/m_Z^2$ 的一阶变分直接给
\begin{align*}
 2s_W\delta s_W
 &= -\delta\!\left(\frac{m_W^2}{m_Z^2}\right)
 =-c_W^2\left(
 \frac{\delta m_W^2}{m_W^2}-\frac{\delta m_Z^2}{m_Z^2}
 \right),
\end{align*}
因此
\begin{equation*}
 \frac{\delta s_W}{s_W}
 =-\frac{c_W^2}{2s_W^2}
 \left(
 \frac{\delta m_W^2}{m_W^2}-\frac{\delta m_Z^2}{m_Z^2}
 \right),
\end{equation*}
得到正文\eqnrefs{eq:sw-counterterm-massrelation-dp11}。
\end{derivation}

\begin{derivation}[从\eqnrefs{eq:delta-rho-definition-dp11}到\eqnrefs{eq:delta-rho-doublet-dp11}]
记 $F_L=(f_1,f_2)_L^T$、$M_i=m_ic^2$。$\Delta\rho$ 只保留破坏守护 $SU(2)$ 的零动量自能差，因此先抽出共同的 $SU(2)_L$ 耦合并定义
\begin{equation*}
\mathcal I(a,b)
\equiv\mu^{4-d}\!\int\frac{\dd^d\ell_E}{(2\pi)^d}
\frac{\ell_E^2}{(\ell_E^2+a)(\ell_E^2+b)},
\qquad a=M_1^2,\quad b=M_2^2,
\end{equation*}
其中已经 Wick 旋转到欧氏能量动量，$\mu$ 是维数正规化尺度。这个标量积分来自左手流的 Dirac 迹。以异味带电流为例，Minkowski 分子为
\begin{align*}
T^{\mu\nu}_{12}
&=\Tr\!\left[\gamma^\mu P_L(\slashed\ell+M_2)
\gamma^\nu P_L(\slashed\ell+M_1)\right]\\
&=\frac12\Tr(\gamma^\mu\slashed\ell\gamma^\nu\slashed\ell)
=2\left(2\ell^\mu\ell^\nu-g^{\mu\nu}\ell^2\right).
\end{align*}
第二步用了 $P_L\gamma^\mu=\gamma^\mu P_R$：两个质量项之间含 $P_RP_L=0$，奇数个 $\gamma$ 的交叉项迹也为零。Lorentz 对称积分
\begin{equation*}
\ell^\mu\ell^\nu\longrightarrow\frac{g^{\mu\nu}}{d}\ell^2
\end{equation*}
把横向自能的零动量系数约化到 $\mathcal I$。暂时抽去两种自能共有的规范耦合、颜色多重度和外部归一化后，带电流的 $(1,2)$ 异味环与第三弱同位旋流的两个同味环组合成唯一的标量差
\begin{equation*}
\mathcal J(a,b)
\equiv
\mathcal I(a,b)-\frac12\mathcal I(a,a)-\frac12\mathcal I(b,b).
\end{equation*}
当 $a=b$ 时三项逐项给出 $\mathcal J(a,a)=0$；同时，大 $\ell_E$ 展开中的共同二次和对数紫外项在这个差中先行相消。

对标量积分作部分分式分解：

\begin{equation*}
\frac{\ell_E^2}{(\ell_E^2+a)(\ell_E^2+b)}
=-\frac{a}{b-a}\frac1{\ell_E^2+a}
+\frac{b}{b-a}\frac1{\ell_E^2+b}
\end{equation*}
并记
\begin{equation*}
A(a)\equiv\mu^{4-d}\int\frac{\dd^d\ell_E}{(2\pi)^d}\frac1{\ell_E^2+a}.
\end{equation*}
于是
\begin{equation*}
\mathcal I(a,b)=\frac{-aA(a)+bA(b)}{b-a},
\qquad
\mathcal I(a,a)=A(a)+aA'(a).
\end{equation*}
在 $d=4-2\epsilon$ 中，$A(a)$ 的展开只需要保留统一的极点/常数部分 $K_\epsilon$ 与质量对数：
\begin{equation*}
A(a)=\frac{a}{16\pi^2}
\left[K_\epsilon+1-\ln\frac{a}{\mu_E^2}\right]+O(\epsilon),
\end{equation*}
这里 $\mu_E$ 是与质量能量同量纲的重整化尺度。代入上述三项组合，$K_\epsilon$、$\ln\mu_E^2$ 和所有与 $a+b$ 成共同局域多项式的部分逐项消失，得到有限结果
\begin{align*}
\mathcal I(a,b)-\frac12\mathcal I(a,a)-\frac12\mathcal I(b,b)
&=\frac{1}{32\pi^2}
\frac{a^2-b^2-2ab\ln(a/b)}{a-b}\\
&=\frac{1}{32\pi^2}
\left[a+b-\frac{2ab}{a-b}\ln\frac ab\right].
\end{align*}

最后恢复弱耦合。带电流顶角给 $g/\sqrt2$，再除以 $m_W^2=g^2v_E^2/(4c^4)$；用能量真空标度 $v_E=v c^2$ 可把共同系数写成 $2/v_E^2$。因此
\begin{align*}
\Delta\rho_f
&=\frac{2N_c}{v_E^2}\mathcal J(a,b)\\
&=\frac{N_c}{16\pi^2v_E^2}
\left[a+b-\frac{2ab}{a-b}\ln\frac ab\right].
\end{align*}
由正文低能匹配 $G_F^{(E)}/\sqrt2=1/(2v_E^2)$，即
\begin{equation*}
\frac1{v_E^2}=\sqrt2\,G_F^{(E)},
\end{equation*}
于是
\begin{equation*}
\Delta\rho_f
=\frac{N_cG_F^{(E)}}{8\pi^2\sqrt2}
\left[
M_1^2+M_2^2-
\frac{2M_1^2M_2^2}{M_1^2-M_2^2}
\ln\frac{M_1^2}{M_2^2}
\right],
\end{equation*}
得到正文\eqnrefs{eq:delta-rho-doublet-dp11}。
\end{derivation}

\begin{derivation}[从\eqnrefs{eq:sm-higgs-cartesian-rxi-dp11}到\eqnrefs{eq:sm-rxi-quadratic-before-gf-dp11}]
要得到 \(R_\xi\) 规范固定前必须消去的混合项，不能只读最终质量矩阵；需要把协变导数在真空附近按场次数展开。正文采用
\begin{equation*}
\Phi
=\sqrt{\hbar c}
\begin{pmatrix}
\widehat\phi^+\\[2pt]
\bigl(v+\widehat h+\ii\widehat\chi\bigr)/\sqrt2
\end{pmatrix}.
\end{equation*}
只保留 \(D_\mu\Phi\) 中对涨落场至多一次的部分。带电上分量可写成
\begin{equation*}
(D_\mu\Phi)_1^{(1)}
=\sqrt{\hbar c}\left(
\partial_\mu\widehat\phi^+
+\ii\kappa_W\mathcal W_\mu^+
\right),
\end{equation*}
其中 \(\kappa_W=m_Wc/\hbar\) 按正文定义。其模平方给
\begin{align*}
\frac{|(D_\mu\Phi)_1^{(1)}|^2}{\hbar c}
={}&\partial_\mu\widehat\phi^-\partial^\mu\widehat\phi^+
+\kappa_W^2\mathcal W^-_\mu\mathcal W^{+\mu}\\
&+\ii\kappa_W
\left(
\mathcal W^-_\mu\partial^\mu\widehat\phi^+
-\mathcal W^+_\mu\partial^\mu\widehat\phi^-
\right).
\end{align*}
因此带电 Goldstone 的动能、\(W\) 的质量二次型以及 \(W\)--Goldstone 导数混合来自同一个平方，三者并不是三个独立假设。

中性下分量的一阶部分同理为
\begin{equation*}
(D_\mu\Phi)_2^{(1)}
=\sqrt{\frac{\hbar c}{2}}
\left[
\partial_\mu\widehat h
+\ii\partial_\mu\widehat\chi
+\ii\kappa_Z\mathcal Z_\mu
\right],
\end{equation*}
其中真空方向对光子组合的系数为零，因此 \(A_\mu^{\rm geo}\) 不产生质量项。取模平方并只保留二次项，实部与虚部分离后得到
\begin{equation*}
\frac12(\partial\widehat h)^2
+\frac12(\partial\widehat\chi)^2
+\frac12\kappa_Z^2\mathcal Z_\mu\mathcal Z^\mu
+\kappa_Z\mathcal Z_\mu\partial^\mu\widehat\chi.
\end{equation*}
这一步同时说明为什么 \(Z\)--Goldstone 混合只涉及 \(\widehat\chi\)，而径向模 \(\widehat h\) 不与 \(Z_\mu\) 发生双线性导数混合。

最后处理势能。若真空由 \(V'(v)=0\) 固定，则在径向方向 Taylor 展开
\begin{equation*}
V(v+\widehat h)
=V(v)+\frac12V''(v)\widehat h^2
+\mathcal O(\widehat h^3),
\end{equation*}
一次项因驻点条件消失。用正文定义 \(\kappa_h^2\equiv V''(v)/(\hbar c)\) 后，拉氏量中的势能贡献为
\begin{equation*}
-\frac12\kappa_h^2\widehat h^2.
\end{equation*}
把带电、中性与径向三部分相加，就逐项得到正文\eqnrefs{eq:sm-rxi-quadratic-before-gf-dp11}。该二次核明确暴露出随后选择 \(R_\xi\) 规范固定函数所要消去的两类导数混合项。
\end{derivation}

\begin{derivation}[从\eqnrefs{eq:sm-rxi-gauge-functions-dp11}到\eqnrefs{eq:sm-rxi-goldstone-masses-dp11}]
先展开带电规范固定项的交叉部分：

\begin{align*}
 -\frac1{\xi_W}F^+F^-
 \supset{}&-\ii\kappa_W\widehat\phi^-\partial\!\cdot\mathcal W^+
 +\ii\kappa_W\widehat\phi^+\partial\!\cdot\mathcal W^-.
\end{align*}
对作用量积分分部，并假设场在积分边界足够快衰减，使表面项为零：
\begin{align*}
 \int\dd^4 x\,
 \widehat\phi^-\partial_\mu\mathcal W^{+\mu}
 &=-\int\dd^4 x\,
 (\partial_\mu\widehat\phi^-)\mathcal W^{+\mu},\\
 \int\dd^4 x\,
 \widehat\phi^+\partial_\mu\mathcal W^{-\mu}
 &=-\int\dd^4 x\,
 (\partial_\mu\widehat\phi^+)\mathcal W^{-\mu}.
\end{align*}
因此规范固定项贡献
\begin{equation*}
 +\ii\kappa_W\mathcal W^+\!\cdot\partial\widehat\phi^-
 -\ii\kappa_W\mathcal W^-\!\cdot\partial\widehat\phi^+,
\end{equation*}
恰与\eqnrefs{eq:sm-rxi-quadratic-before-gf-dp11} 的带电混合项逐项相消。中性部分为
\begin{equation*}
 -\frac1{2\xi_Z}(F^Z)^2
 \supset
 +\kappa_Z\widehat\chi\,\partial\!\cdot\mathcal Z
 \longrightarrow
 -\kappa_Z\mathcal Z\!\cdot\partial\widehat\chi,
\end{equation*}
同样消去 $Z$--$\chi$ 混合。

规范固定函数平方还留下纯标量项
\begin{equation*}
 -\xi_W\kappa_W^2\widehat\phi^+\widehat\phi^-
 -\frac12\xi_Z\kappa_Z^2\widehat\chi^2.
\end{equation*}
与标量传播子的标准二次质量项比较，得到
\begin{equation*}
 m_{\phi^\pm}^2=\xi_Wm_W^2,
 \qquad
 m_\chi^2=\xi_Zm_Z^2,
\end{equation*}
即\eqnrefs{eq:sm-rxi-goldstone-masses-dp11}。
\end{derivation}

\begin{derivation}[从\eqnrefs{eq:sm-rxi-gauge-functions-dp11}到\eqnrefs{eq:sm-rxi-ghost-masses-dp11}]
由正文的 $\mathcal M_{\rm FP}^{ab}(x,y)=\delta F^a(x)/\delta\alpha^b(y)$ 出发。在线性化真空附近，破缺规范变换使 Goldstone 场沿规范轨道移动：

\begin{equation*}
 \delta\widehat\phi^\pm
 =\pm\ii\kappa_W\alpha^\pm+\order(\text{场}),
 \qquad
 \delta\widehat\chi
 =\kappa_Z\alpha^Z+\order(\text{场}).
\end{equation*}
同时规范场的线性变换给
\begin{equation*}
 \delta(\partial\!\cdot\mathcal W^\pm)=\partial^2\alpha^\pm+\order(\text{场}),
 \qquad
 \delta(\partial\!\cdot\mathcal Z)=\partial^2\alpha^Z+\order(\text{场}),
\end{equation*}
以及 $\delta(\partial\!\cdot A^{\rm geo})=\partial^2\alpha^\gamma+\order(\text{场})$。代入\eqnrefs{eq:sm-rxi-gauge-functions-dp11}：
\begin{align*}
 \delta F^\pm
 &=\left(\partial^2+\xi_W\kappa_W^2\right)\alpha^\pm
 +\order(\text{场}),\\
 \delta F^Z
 &=\left(\partial^2+\xi_Z\kappa_Z^2\right)\alpha^Z
 +\order(\text{场}),\\
 \delta F^\gamma
 &=\partial^2\alpha^\gamma+\order(\text{场}).
\end{align*}
于是鬼场二次算符分别是
\begin{equation*}
 \partial^2+\xi_W\kappa_W^2,
 \qquad
 \partial^2+\xi_Z\kappa_Z^2,
 \qquad
 \partial^2,
\end{equation*}
与自由标量二次核比较便得到\eqnrefs{eq:sm-rxi-ghost-masses-dp11}。
\end{derivation}

\begin{derivation}[从\eqnrefs{eq:sm-rxi-gf-lagrangian-dp11}到\eqnrefs{eq:sm-rxi-massive-vector-propagator-dp41}]
取物理四动量 $p^\mu$，并记

\begin{equation*}
 P^T_{\mu\nu}=\eta_{\mu\nu}-\frac{p_\mu p_\nu}{p^2},
 \qquad
 P^L_{\mu\nu}=\frac{p_\mu p_\nu}{p^2}.
\end{equation*}
它们满足
\begin{equation*}
 (P^T)^2=P^T,
 \qquad
 (P^L)^2=P^L,
 \qquad
 P^TP^L=0,
 \qquad
 P^T+P^L=I.
\end{equation*}
规范固定后的约化二次核可写成
\begin{equation*}
 K_{\mu\nu}(p)
 =-(p^2-m_V^2c^2)\eta_{\mu\nu}
 +\left(1-\frac1{\xi_V}\right)p_\mu p_\nu.
\end{equation*}
用 $\eta=P^T+P^L$ 与 $p_\mu p_\nu=p^2P^L_{\mu\nu}$ 分解：
\begin{align*}
 K_{\mu\nu}
 ={}&-(p^2-m_V^2c^2)P^T_{\mu\nu}\\
 &-\frac1{\xi_V}(p^2-\xi_Vm_V^2c^2)P^L_{\mu\nu}.
\end{align*}
两个投影子空间互相正交，所以逆核逐块取倒数：
\begin{equation*}
 K^{-1}_{\mu\nu}
 =-\frac{P^T_{\mu\nu}}{p^2-m_V^2c^2}
 -\frac{\xi_VP^L_{\mu\nu}}{p^2-\xi_Vm_V^2c^2}.
\end{equation*}
乘上 Feynman 边界条件对应的 $\ii$ 并重新合并两个投影算符，得到
\begin{equation*}
 D^{(V)}_{\mu\nu}(p)
 =\frac{-\ii}{p^2-m_V^2c^2+\ii0^+}
 \left[
 \eta_{\mu\nu}
 -(1-\xi_V)
 \frac{p_\mu p_\nu}{p^2-\xi_Vm_V^2c^2+\ii0^+}
 \right],
\end{equation*}
即\eqnrefs{eq:sm-rxi-massive-vector-propagator-dp41}。$\xi_V=1$ 时纵向附加项消失；$\xi_V\to\infty$ 时传播子化为幺正规范形式。这两个极限分别检验了纵向极点结构与质量项的归一化。
\end{derivation}

\begin{derivation}[从\eqnrefs{eq:sm-higgs-polar-scattering-dp87}到\eqnrefs{eq:ww-high-energy-constant-dp11}]
把
\[
 \pi\equiv\pi^a\sigma^a,
 \qquad
 U_\pi=\ee^{\ii\pi/v}
\]
展开到产生四 Goldstone 顶角所需的四阶。指数展开为
\begin{align*}
 U_\pi
 &=1+\frac{\ii\pi}{v}-\frac{\pi^2}{2v^2}
 -\frac{\ii\pi^3}{6v^3}+\frac{\pi^4}{24v^4}+\order(\pi^5),\\
 U_\pi^\dagger
 &=1-\frac{\ii\pi}{v}-\frac{\pi^2}{2v^2}
 +\frac{\ii\pi^3}{6v^3}+\frac{\pi^4}{24v^4}+\order(\pi^5).
\end{align*}
利用 $\Tr(\sigma^a\sigma^b)=2\delta^{ab}$ 和
\[
 \sigma^a\sigma^b=\delta^{ab}I+\ii\epsilon^{abc}\sigma^c,
\]
把正文\eqnrefs{eq:sm-goldstone-kinetic-polar-dp87} 展开到 $\pi^4$，得到规范归一化的 Goldstone 动能以及第一项非线性相互作用
\begin{equation*}
 \Lag_{\pi}
 =\frac12\partial_\mu\pi^a\partial^\mu\pi^a
 +\frac{1}{6v^2}
 \left[
 (\pi^a\partial_\mu\pi^a)^2
 -\pi^2(\partial_\mu\pi^b\partial^\mu\pi^b)
 \right]
 +\order(\pi^6/v^4).
\end{equation*}
同时 $(v+h)^2/v^2=1+2h/v+h^2/v^2$ 给出
\begin{equation*}
 \Lag_{h\pi\pi}
 =\frac{h}{v}\,
 \partial_\mu\pi^a\partial^\mu\pi^a+\frac{h^2}{2v^2}
 \partial_\mu\pi^a\partial^\mu\pi^a+\cdots .
\end{equation*}

定义带电与中性组合
\begin{equation*}
 w^\pm=\frac{\pi^1\mp\ii\pi^2}{\sqrt2},
 \qquad z=\pi^3.
\end{equation*}
对 $w^+(p_1)w^-(p_2)\to z(p_3)z(p_4)$，所有外腿在高能极限可取无质量，故
\begin{equation*}
 p_i^2=0,
 \qquad
 c^2s=c^2(p_1+p_2)^2=2c^2p_1\cdot p_2=2c^2p_3\cdot p_4.
\end{equation*}
从四 Goldstone 导数项逐个把导数替换为外腿动量，四点顶角给
\begin{equation*}
 \mathcal M_{4\pi}
 =\frac{c^2s}{v_E^2},
 \qquad v_E^2=(\hbar c)v^2.
\end{equation*}
这里 $1/v^2$ 与量子化场的 SI 归一化合并为 $c^2/v_E^2$；这正是没有可解析径向模时的非线性 Goldstone 振幅。

再计算 $s$ 道 Higgs 交换。由 $h\pi\pi$ 导数顶角，两个在壳顶角分别给
\begin{equation*}
 h w^+w^-:\quad
 \ii\frac{c^2s}{v_E},
 \qquad
 hzz:\quad
 \ii\frac{c^2s}{v_E},
\end{equation*}
而径向模传播子为
\begin{equation*}
 \frac{\ii}{c^2s-(m_hc^2)^2+\ii0}.
\end{equation*}
因此
\begin{equation*}
 \mathcal M_h
 =-\frac{c^4s^2}{v_E^2[c^2s-(m_hc^2)^2]}.
\end{equation*}
把共同的 $c^2$ 从分母约去，等价写成正文使用的物理四动量 convention：
\begin{align*}
 \mathcal M
 &=\mathcal M_{4\pi}+\mathcal M_h\\
 &=\frac{c^2s}{v_E^2}
 -\frac{c^2s^2}{v_E^2(s-m_h^2c^2)}\\
 &=\frac{c^2s(s-m_h^2c^2)-c^2s^2}
 {v_E^2(s-m_h^2c^2)}\\
 &=-\frac{(m_hc^2)^2}{v_E^2}
 \frac{s}{s-m_h^2c^2},
\end{align*}
得到正文\eqnrefs{eq:ww-higgs-full-amplitude-dp32}。现在高能抵消不是从最后两项“看出来”的，而是由同一母动能项产生的四点接触顶角和 $h\pi\pi$ 顶角共同保证。

最后令 $s\gg m_h^2c^2$，作受控展开
\begin{equation*}
 \frac{s}{s-m_h^2c^2}
 =\frac1{1-m_h^2c^2/s}
 =1+\frac{m_h^2c^2}{s}
 +\order\!\left(\frac{m_h^4c^4}{s^2}\right),
\end{equation*}
于是
\begin{equation*}
 \mathcal M
 =-\frac{(m_hc^2)^2}{v_E^2}
 \left[1+\order\!\left(\frac{m_h^2c^2}{s}\right)\right],
\end{equation*}
其首项即正文\eqnrefs{eq:ww-high-energy-constant-dp11}。高能展开的两个独立控制量分别是 Goldstone 等价定理的 $m_Wc^2/E$ 和此处径向传播子的 $m_h^2c^2/s$。
\end{derivation}

\begin{derivation}[从\eqnrefs{eq:sm-coleman-weinberg-si-dp12}到\eqnrefs{eq:sm-higgs-quartic-beta-dp93}]
四次耦合的 beta 函数由有效势中 $\varphi_E^4$ 的紫外对数系数决定。对任意关于背景场的多项式 $P(\varphi_E)$，定义系数提取算符
\begin{equation*}
 [\varphi_E^4]P
 \equiv\frac1{4!}\left.\frac{\dd^4P}{\dd\varphi_E^4}\right|_{\varphi_E=0}.
\end{equation*}
一圈背景质量一般可写成
\begin{equation*}
 \mathcal M_i^2(\varphi_E)=m_{i,0}^2+\kappa_i\varphi_E^2,
\end{equation*}
于是无需取大背景极限便有精确代数恒等式
\begin{equation*}
 [\varphi_E^4]\,\mathcal M_i^4(\varphi_E)=\kappa_i^2.
\end{equation*}
正文\eqnrefs{eq:sm-coleman-weinberg-si-dp12} 对重整化能标的显式导数为
\begin{equation*}
 E_{\rm R}\frac{\partial V_1}{\partial E_{\rm R}}
 =-\frac1{32\pi^2(\hbar c)^3}
 \sum_i n_i\mathcal M_i^4(\varphi_E),
\end{equation*}
所以其四次背景场系数严格为
\begin{equation*}
 [\varphi_E^4]\left(E_{\rm R}\frac{\partial V_1}{\partial E_{\rm R}}\right)
 =-\frac1{32\pi^2(\hbar c)^3}\sum_i n_\ii\kappa_i^2.
\end{equation*}
逐类计算无量纲系数 $\sum_i n_\ii\kappa_i^2$。

先算标量涨落。由树级势的 Hessian，径向 Higgs 与三个 Goldstone 模的 $\varphi_E^2$ 系数满足
\begin{equation*}
 \kappa_h=3\lambda_H,
 \qquad
 \kappa_G=\lambda_H,
 \qquad
 n_h=1,
 \qquad
 n_G=3,
\end{equation*}
因此
\begin{equation*}
 \left.\sum_i n_\ii\kappa_i^2\right|_{\rm scalar}
 =(3\lambda_H)^2+3\lambda_H^2
 =12\lambda_H^2.
\end{equation*}

再算规范涨落。$W^\pm$ 共给六个有质量物理极化自由度，$Z$ 给三个，而
\begin{equation*}
 \kappa_W=\frac{g^2}{4},
 \qquad
 \kappa_Z=\frac{g^2+g'^2}{4}.
\end{equation*}
于是
\begin{align*}
 \left.\sum_i n_\ii\kappa_i^2\right|_{\rm gauge}
 &=6\left(\frac{g^2}{4}\right)^2
 +3\left(\frac{g^2+g'^2}{4}\right)^2\\
 &=\frac38g^4+\frac{3}{16}(g^2+g'^2)^2.
\end{align*}

最后处理全部 Yukawa 矩阵，而不是预先只留下顶夸克。对任意一个 Dirac 质量本征值 $y_a$，背景质量满足
\begin{equation*}
 \mathcal M_a^2=\frac{y_a^2}{2}\varphi_E^2,
\end{equation*}
一个 Dirac 费米子给四个自旋/反粒子 Grassmann 自由度，故 $n_a=-4N_c$。其贡献为
\begin{equation*}
 n_a\kappa_a^2
 =-4N_c\left(\frac{y_a^2}{2}\right)^2
 =-N_cy_a^4.
\end{equation*}
对所有上型夸克、下型夸克和带电轻子求和，并把质量本征值平方的和改写成味基不变量，得到
\begin{align*}
 \left.\sum_i n_\ii\kappa_i^2\right|_{\rm Yukawa}
 &=-\Tr\!\left[3(Y_uY_u^\dagger)^2
 +3(Y_dY_d^\dagger)^2
 +(Y_eY_e^\dagger)^2\right]\\
 &=-\mathcal H_Y.
\end{align*}
因此全部一圈行列式的显式尺度导数由
\begin{equation*}
 \mathcal S_4
 \equiv\sum_i n_\ii\kappa_i^2
 =12\lambda_H^2
 +\frac38g^4
 +\frac{3}{16}(g^2+g'^2)^2
 -\mathcal H_Y
\end{equation*}
控制。

仅有显式对数还不够，因为正文\eqnrefs{eq:sm-effective-potential-rg-equation-dp12} 同时包含背景 Higgs 场的反常维数。在同一 $\overline{\rm MS}$ 一圈阶，最小标准模型 Higgs 双重态满足
\begin{equation*}
 16\pi^2\gamma_\Phi
 =\mathcal T_Y-\frac94g^2-\frac34g'^2.
\end{equation*}
树级四次势为
\begin{equation*}
 V_0^{(4)}
 =\frac{\lambda_H\varphi_E^4}{4(\hbar c)^3}.
\end{equation*}
把它与一圈显式尺度导数一起代入 Callan--Symanzik 方程，并只保留 $\varphi_E^4$ 与一圈量级：
\begin{equation*}
 0
 =-\frac{\mathcal S_4}{32\pi^2}
 +\frac{\beta_{\lambda_H}}4
 -\gamma_\Phi\lambda_H.
\end{equation*}
解出四次耦合的 beta 函数，
\begin{equation*}
 \beta_{\lambda_H}
 =\frac{\mathcal S_4}{8\pi^2}
 +4\gamma_\Phi\lambda_H.
\end{equation*}
乘以 $16\pi^2$ 并代入 $\mathcal S_4$ 与 $\gamma_\Phi$：
\begin{align*}
16\pi^2\beta_{\lambda_H}
={}&24\lambda_H^2
+\frac34g^4
+\frac38(g^2+g'^2)^2
-2\mathcal H_Y\\
&+4\lambda_H\mathcal T_Y
-9g^2\lambda_H-3g'^2\lambda_H\\
={}&24\lambda_H^2
-3\lambda_H(3g^2+g'^2)
+\frac38\left[2g^4+(g^2+g'^2)^2\right]
+4\lambda_H\mathcal T_Y-2\mathcal H_Y,
\end{align*}
即正文\eqnrefs{eq:sm-higgs-quartic-beta-dp93}。Higgs 双重态是 $SU(3)_c$ 单态，因此这一圈背景质量谱中没有由 Higgs 背景直接产生的胶子质量项，$\beta_{\lambda_H}$ 在该圈阶没有显式 $g_s$ 项；强耦合通过 $Y_u,Y_d$ 的重整化群方程间接进入耦合的 RG 系统。
\end{derivation}

\begin{derivation}[从\eqnrefs{eq:sm-higgs-lagrangian-dp11}到\eqnrefs{eq:sm-higgs-mass-dp11}]
把正文 Higgs 势写成单一规范不变量
\begin{equation*}
 x\equiv\Phi^\dagger\Phi\ge0,
 \qquad
 V(x)=-\mu_H^2x+\frac{\lambda_H}{\hbar c}x^2.
\end{equation*}
对复二重态作任意变分时，$\delta x=\delta\Phi^\dagger\Phi+\Phi^\dagger\delta\Phi$，因此驻点满足
\begin{equation*}
 \frac{\partial V}{\partial\Phi^\dagger}
 =\left(-\mu_H^2+\frac{2\lambda_H}{\hbar c}\Phi^\dagger\Phi\right)\Phi=0.
\end{equation*}
除原点 $\Phi=0$ 外，非零驻点组成
\begin{equation*}
 \Phi^\dagger\Phi
 =\frac{\mu_H^2\hbar c}{2\lambda_H}.
\end{equation*}
当 $\mu_H^2>0$、$\lambda_H>0$ 时，原点沿径向的二阶变化为负，而上述非零流形是稳定极小值。利用电弱规范变换可以从这条真空轨道中选取代表元
\begin{equation*}
 \langle\Phi\rangle
 =\frac1{\sqrt2}\binom0v,
\end{equation*}
于是
\begin{equation*}
 \frac{v^2}{2}
 =\frac{\mu_H^2\hbar c}{2\lambda_H},
 \qquad
 v^2=\frac{\mu_H^2\hbar c}{\lambda_H},
 \qquad
 v_E\equiv\sqrt{\hbar c}\,v,
\end{equation*}
即正文\eqnrefs{eq:sm-higgs-vev-dp11}。

真空附近先保留一般规范轨道方向，写成
\begin{equation*}
 \Phi(x)=\frac1{\sqrt2}U_\xi(x)\binom0{v+h(x)}.
\end{equation*}
势能只依赖 $\Phi^\dagger\Phi$，所以 $U_\xi^\dagger U_\xi=I$ 使三个切向场 $\xi^I$ 完全不进入 $V$；势能只约束径向变量 $h$。取幺正规范 $U_\xi=I$ 后，
\begin{equation*}
 \Phi^\dagger\Phi=\frac12(v+h)^2,
\end{equation*}
代回势能：
\begin{align*}
 V(h)
 &=-\frac{\mu_H^2}{2}(v+h)^2
 +\frac{\lambda_H}{4\hbar c}(v+h)^4\\
 &=-\frac{\mu_H^2}{2}
 (v^2+2vh+h^2)
 +\frac{\lambda_H}{4\hbar c}
 (v^4+4v^3h+6v^2h^2+4vh^3+h^4).
\end{align*}
逐阶收集：
\begin{align*}
 V(h)=V_0
 &+v\left(-\mu_H^2+\frac{\lambda_Hv^2}{\hbar c}\right)h\\
 &+\left(-\frac{\mu_H^2}{2}+\frac{3\lambda_Hv^2}{2\hbar c}\right)h^2
 +\frac{\lambda_Hv}{\hbar c}h^3
 +\frac{\lambda_H}{4\hbar c}h^4.
\end{align*}
利用真空条件 $\lambda_Hv^2/(\hbar c)=\mu_H^2$，线性项消失，二次项化为
\begin{equation*}
 V(h)=V_0+\mu_H^2h^2
 +\frac{\lambda_Hv}{\hbar c}h^3
 +\frac{\lambda_H}{4\hbar c}h^4,
\end{equation*}
即正文\eqnrefs{eq:sm-higgs-potential-expanded-dp11}。

最后核对径向场的动能归一化。忽略规范场时，
\begin{equation*}
 (\partial_\mu\Phi_U)^\dagger(\partial^\mu\Phi_U)
 =\frac12\partial_\mu h\,\partial^\mu h,
\end{equation*}
所以 $h$ 已经是标准规范化的实标量。标准质量项为
\begin{equation*}
 -\frac12\left(\frac{m_hc}{\hbar}\right)^2h^2.
\end{equation*}
由于拉格朗日量含 $-V$，比较二次项得到
\begin{equation*}
 \frac12\left(\frac{m_hc}{\hbar}\right)^2=\mu_H^2.
\end{equation*}
再用
\begin{equation*}
 \mu_H^2=\frac{\lambda_Hv^2}{\hbar c}
 =\frac{\lambda_Hv_E^2}{(\hbar c)^2},
\end{equation*}
便有
\begin{equation*}
 (m_hc^2)^2=2\lambda_Hv_E^2,
 \qquad
 \boxed{m_hc^2=\sqrt{2\lambda_H}\,v_E},
\end{equation*}
即正文\eqnrefs{eq:sm-higgs-mass-dp11}。同一展开同时固定 $h^3$ 和 $h^4$ 自耦，因此质量、三点自耦和四点自耦都来自同一个径向势，而不是彼此独立的参数。
\end{derivation}





\begin{derivation}[从\eqnrefs{eq:muon-full-w-amplitude-dp11}到\eqnrefs{eq:fermi-amplitude-pl-dp11}]
正文\eqnrefs{eq:muon-full-w-amplitude-dp11} 仍保留中间 $W$ 的完整动量依赖。实验中的 $\mu$ 子衰变满足 $|q^2|\ll m_W^2c^2$。在这一控制参数下展开完整 $W$ 传播子，局域四费米相互作用出现在最低阶。

传播子分母的几何展开为
\begin{align*}
\frac{1}{q^2-m_W^2c^2}
&=-\frac{1}{m_W^2c^2}\frac{1}{1-q^2/(m_W^2c^2)}\\
&=-\frac{1}{m_W^2c^2}
\left[1+\frac{q^2}{m_W^2c^2}
+\mathcal O\!\left(\frac{q^4}{m_W^4c^4}\right)\right].
\end{align*}
纵向分子与轻费米子流收缩时带有外腿质量。以 $\mu$ 子流为例，$q=p-k$，
\begin{align*}
q_\alpha\bar u_{\nu_\mu}(k)\gamma^\alpha P_Lu_\mu(p)
&=\bar u_{\nu_\mu}(k)(\slashed p-\slashed k)P_Lu_\mu(p)\\
&=m_\mu c\,\bar u_{\nu_\mu}P_Ru_\mu
-m_{\nu_\mu}c\,\bar u_{\nu_\mu}P_Lu_\mu.
\end{align*}
电子一侧同理得到
\begin{equation*}
q_\beta\bar u_e\gamma^\beta P_Lv_{\bar\nu_e}
=m_{\mathrm e} c\,\bar u_eP_Lv_{\bar\nu_e}
-m_{\nu_e}c\,\bar u_eP_Rv_{\bar\nu_e}.
\end{equation*}
因此纵向部分至少带两个轻费米子质量插入；相对于横向电流--电流项，它在振幅中受 $m_fm_{f'}/m_W^2$ 型比值抑制。保留最低阶 $\eta_{\alpha\beta}$ 后，先只看两个顶角与传播子的标量系数：
\begin{align*}
\left(-\frac{\ii g}{\sqrt2}\right)^2
\frac{-\ii}{q^2-m_W^2c^2}
&=\frac{\ii g^2}{2(q^2-m_W^2c^2)}\\
&=-\ii\frac{g^2}{2m_W^2c^2}
\left[1+\frac{q^2}{m_W^2c^2}
+\mathcal O\!\left(\frac{q^4}{m_W^4c^4}\right)\right].
\end{align*}
于是完整振幅在最低阶成为
\begin{align*}
\mathcal M
={}&-\frac{g^2}{2m_W^2c^2}
\bigl[\bar u_{\nu_\mu}\gamma^\alpha P_Lu_\mu\bigr]
\bigl[\bar u_e\gamma_\alpha P_Lv_{\bar\nu_e}\bigr]\\
&+\mathcal O\!\left(
\frac{q^2}{m_W^2c^2},
\frac{m_fm_{f'}}{m_W^2}
\right).
\end{align*}
这正是正文\eqnrefs{eq:fermi-amplitude-pl-dp11}。
\end{derivation}

\begin{derivation}[从\eqnrefs{eq:fermi-amplitude-pl-dp11}到\eqnrefs{eq:muon-lifetime-tree-dp11}]
正文\eqnrefs{eq:fermi-amplitude-pl-dp11} 已把虚 $W$ 消去，只剩两条左手带电流。自旋求和由 Dirac 完备关系计算，三体相空间则分解成连续两次 Lorentz 协变二体相空间积分。

为避免在相空间中反复携带 $c$，把四动量改写为能量量纲
\begin{equation*}
P\equiv cp,\qquad P'\equiv cp',\qquad K\equiv ck,\qquad K'\equiv ck',
\end{equation*}
并记 $M\equiv m_\mu c^2$。利用 $P_L=(1-\gamma^5)/2$ 与正文\eqnrefs{eq:gf-matching-dp11}，四费米振幅可等价写成通常的 $V-A$ 归一化
\begin{equation*}
\mathcal M
=-\frac{G_F^{(E)}}{\sqrt2}
\bigl[\bar u(K)\gamma^\alpha(1-\gamma^5)u(P)\bigr]
\bigl[\bar u(P')\gamma_\alpha(1-\gamma^5)v(K')\bigr].
\end{equation*}
忽略电子和中微子质量，并对初态 $\mu$ 子两个自旋取平均。由
\[
\sum_su(P,s)\bar u(P,s)=\slashed P+M,
\qquad
\sum_su(K,s)\bar u(K,s)=\slashed K,
\qquad
\sum_sv(K',s)\bar v(K',s)=\slashed K'
\]
得到
\begin{align*}
\overline{|\mathcal M|^2}
&=\frac{[G_F^{(E)}]^2}{4}\,
L^{\alpha\beta}H_{\alpha\beta},\\
L^{\alpha\beta}
&\equiv
\Tr\!\left[
\slashed K\gamma^\alpha(1-\gamma^5)
(\slashed P+M)\gamma^\beta(1-\gamma^5)
\right],\\
H_{\alpha\beta}
&\equiv
\Tr\!\left[
\slashed P'\gamma_\alpha(1-\gamma^5)
\slashed K'\gamma_\beta(1-\gamma^5)
\right].
\end{align*}
这里前因子 $1/4$ 同时包含振幅平方中的 $1/2$ 和初态自旋平均的 $1/2$。

先说明 $L^{\alpha\beta}$ 中的 $M$ 项为什么消失。利用
$(1-\gamma^5)\gamma^\beta=\gamma^\beta(1+\gamma^5)$，
\begin{align*}
&M\Tr\!\left[
\slashed K\gamma^\alpha(1-\gamma^5)
\gamma^\beta(1-\gamma^5)
\right]\\
&\quad=M\Tr\!\left[
\slashed K\gamma^\alpha\gamma^\beta
(1+\gamma^5)(1-\gamma^5)
\right]=0.
\end{align*}
因此只需计算四个普通 $\gamma$ 矩阵的迹与含一个 $\gamma^5$ 的迹。采用
\begin{align*}
\Tr(\gamma^\rho\gamma^\alpha\gamma^\sigma\gamma^\beta)
&=4(\eta^{\rho\alpha}\eta^{\sigma\beta}
-\eta^{\rho\sigma}\eta^{\alpha\beta}
+\eta^{\rho\beta}\eta^{\alpha\sigma}),\\
\Tr(\gamma^\rho\gamma^\alpha\gamma^\sigma\gamma^\beta\gamma^5)
&=-4\ii\epsilon^{\rho\alpha\sigma\beta},
\end{align*}
并用 $(1-\gamma^5)^2=2(1-\gamma^5)$，逐项得到
\begin{align*}
L^{\alpha\beta}
&=8\left(
K^\alpha P^\beta+K^\beta P^\alpha
-\eta^{\alpha\beta}K\!\cdot P
+\ii\epsilon^{\alpha\beta\rho\sigma}K_\rho P_\sigma
\right),\\
H_{\alpha\beta}
&=8\left(
P'_\alpha K'_\beta+P'_\beta K'_\alpha
-\eta_{\alpha\beta}P'\!\cdot K'
+\ii\epsilon_{\alpha\beta\mu\nu}P'^\mu K'^\nu
\right).
\end{align*}
把每个张量分成对称部分 $S$ 与反对称部分 $A$。由于 $S^{\alpha\beta}A_{\alpha\beta}=0$，交叉项消失。对称部分缩并为
\begin{equation*}
S_L^{\alpha\beta}S_{H,\alpha\beta}
=2\left[(K\!\cdot P')(P\!\cdot K')
+(K\!\cdot K')(P\!\cdot P')\right],
\end{equation*}
而反对称部分用
\begin{equation*}
\epsilon^{\alpha\beta\rho\sigma}
\epsilon_{\alpha\beta\mu\nu}
=-2\left(\delta^\rho_{\mu}\delta^\sigma_{\nu}
-\delta^\rho_{\nu}\delta^\sigma_{\mu}\right)
\end{equation*}
得到
\begin{equation*}
A_L^{\alpha\beta}A_{H,\alpha\beta}
=2\left[(K\!\cdot P')(P\!\cdot K')
-(K\!\cdot K')(P\!\cdot P')\right].
\end{equation*}
两者相加后第二种标量积精确相消，于是
\begin{equation*}
L^{\alpha\beta}H_{\alpha\beta}
=256(K\!\cdot P')(P\!\cdot K'),
\end{equation*}
从而
\begin{equation*}
\boxed{
\overline{|\mathcal M|^2}
=64[G_F^{(E)}]^2(P\!\cdot K')(P'\!\cdot K).}
\end{equation*}
关键结构现在显式可见：$V-A$ 的反对称 $\epsilon$ 张量项恰好删除了另一种标量积配对，而不是给平方振幅额外增加一个独立项。

对两个不可观测的中微子积分。定义
\begin{equation*}
Q\equiv K+K'=P-P',\qquad s_\nu\equiv Q^2.
\end{equation*}
三体 Lorentz 不变相空间先插入恒等分解
\[
1=\int\frac{\dd s_\nu}{2\pi}
\int\dd^4Q\,(2\pi)^4
\delta^{(4)}(Q-K-K')\delta(Q^2-s_\nu)\Theta(Q^0),
\]
于是严格因子化为
\begin{equation*}
\dd\mathfrak P_3
=\frac{\dd s_\nu}{2\pi}
\dd\mathfrak P_2(P;P',Q)\,
\dd\mathfrak P_2(Q;K,K').
\end{equation*}
固定 $Q$ 后只需计算二阶张量积分
\begin{equation*}
I^{\alpha\beta}(Q)
\equiv\int\dd\mathfrak P_2(Q;K,K')K^\alpha K'^\beta.
\end{equation*}
积分后只剩 $Q^\mu$ 与 $\eta^{\mu\nu}$，故
\begin{equation*}
I^{\alpha\beta}=A s_\nu\eta^{\alpha\beta}+B Q^\alpha Q^\beta.
\end{equation*}
系数 $A,B$ 不需要猜。到 $Q$ 静止系，令 $\sqrt{s_\nu}>0$，则
\begin{equation*}
K^\mu=\frac{\sqrt{s_\nu}}2(1,\mathbf n),\qquad
K'^\mu=\frac{\sqrt{s_\nu}}2(1,-\mathbf n),\qquad
\dd\mathfrak P_2=\frac{\dd\Omega}{32\pi^2}.
\end{equation*}
因此
\begin{align*}
I^{00}
&=\frac{s_\nu}{4}\frac{4\pi}{32\pi^2}
=\frac{s_\nu}{32\pi},\\
I^{0i}&=0,\\
I^{ij}
&=-\frac{s_\nu}{4}\frac{1}{32\pi^2}
\int\dd\Omega\,n_in_j
=-\frac{s_\nu}{96\pi}\delta_{ij}.
\end{align*}
与协变分解逐分量比较即得
\begin{equation*}
A=\frac1{96\pi},\qquad B=\frac1{48\pi},
\end{equation*}
即
\begin{equation*}
I^{\alpha\beta}(Q)
=\frac{1}{96\pi}
\left(s_\nu\eta^{\alpha\beta}+2Q^\alpha Q^\beta\right).
\end{equation*}
所以
\begin{align*}
&\int\dd\mathfrak P_2(Q;K,K')
(P\!\cdot K')(P'\!\cdot K)\\
&\qquad=P_\alpha P'_\beta I^{\beta\alpha}
=\frac{1}{96\pi}
\left[s_\nu(P\!\cdot P')+2(P\!\cdot Q)(P'\!\cdot Q)\right].
\end{align*}

在 $\mu$ 子静止系 $P=(M,\mathbf0)$，令
$x=2E_e/M$。无质量电子满足 $P'^2=0$，故
\begin{align*}
P\!\cdot P'&=\frac{M^2x}{2},\\
s_\nu=(P-P')^2&=M^2(1-x),\\
P\!\cdot Q&=M^2\left(1-\frac{x}{2}\right),\\
P'\!\cdot Q&=\frac{M^2x}{2}.
\end{align*}
代入可得
\begin{equation*}
\int\dd\mathfrak P_2(Q;K,K')\overline{|\mathcal M|^2}
=\frac{[G_F^{(E)}]^2M^4}{3\pi}x(3-2x).
\end{equation*}
外层二体相空间由一般公式
\begin{equation*}
\dd\mathfrak P_2(P;P',Q)
=\frac{|\mathbf P'|}{16\pi^2M}\dd\Omega_e
\end{equation*}
以及 $|\mathbf P'|=E_e=Mx/2$ 化为
\begin{equation*}
\dd\mathfrak P_2(P;P',Q)
=\frac{x}{32\pi^2}\dd\Omega_e,
\qquad
|\dd s_\nu|=M^2\dd x.
\end{equation*}
于是
\begin{align*}
\frac{\dd\Gamma_\mu^{(0)}}{\dd x}
&=\frac{1}{2M\hbar}\frac1{2\pi}
\left[\frac{[G_F^{(E)}]^2M^4}{3\pi}x(3-2x)\right]
\left[\frac{x}{32\pi^2}\int\dd\Omega_e\right]M^2\\
&=\frac{[G_F^{(E)}]^2M^5}{96\pi^3\hbar}
 x^2(3-2x),
\end{align*}
其中 $\int\dd\Omega_e=4\pi$。这就是正文\eqnrefs{eq:muon-differential-tree-dp73}。

总宽度现在只剩单变量初等积分：
\begin{align*}
\Gamma_\mu^{(0)}
&=\frac{[G_F^{(E)}]^2M^5}{96\pi^3\hbar}
\int_0^1\dd x\,x^2(3-2x)\\
&=\frac{[G_F^{(E)}]^2M^5}{96\pi^3\hbar}
\left[x^3-\frac{x^4}{2}\right]_0^1\\
&=\frac{[G_F^{(E)}]^2(m_\mu c^2)^5}{192\pi^3\hbar},
\end{align*}
即正文\eqnrefs{eq:muon-lifetime-tree-dp11}。
\end{derivation}

\begin{derivation}[从\eqnrefs{eq:generic-vector-ff-exact-width-dp92}到\eqnrefs{eq:zff-width-tree-dp11}]
令
\begin{equation*}
 q_V\equiv\frac{M_V}{c},\qquad a\equiv m_1c,\qquad b\equiv m_2c,
\end{equation*}
于是 $P^2=q_V^2$、$p^2=a^2$、$k^2=b^2$。先把手征顶角改写为矢量--轴矢量形式
\begin{equation*}
 \Gamma^\mu
 =\gamma^\mu(g_V-g_A\gamma^5),\qquad
 g_V=\frac{c_L+c_R}{2},\qquad
 g_A=\frac{c_L-c_R}{2}.
\end{equation*}
对末态自旋求和得到
\begin{align*}
 T^{\mu\nu}
 &=\Tr\!\left[(\slashed p+a)\gamma^\mu(g_V-g_A\gamma^5)
 (\slashed k-b)\gamma^\nu(g_V^*-g_A^*\gamma^5)\right].
\end{align*}
极化投影张量关于 $\mu,\nu$ 对称，因此含 $\epsilon^{\mu\nu\alpha\beta}$ 的矢量--轴矢量交叉项不贡献总宽度。两个对称迹分别为
\begin{align*}
 T_{VV}^{\mu\nu}
 &=4\left[p^\mu k^\nu+p^\nu k^\mu
 -\eta^{\mu\nu}(p\!\cdot\!k+ab)\right],\\
 T_{AA}^{\mu\nu}
 &=4\left[p^\mu k^\nu+p^\nu k^\mu
 -\eta^{\mu\nu}(p\!\cdot\!k-ab)\right].
\end{align*}
利用
\begin{equation*}
 |g_V|^2+|g_A|^2=\frac{|c_L|^2+|c_R|^2}{2},\qquad
 |g_A|^2-|g_V|^2=-\operatorname{Re}(c_Lc_R^*),
\end{equation*}
可把对称部分整理为
\begin{equation*}
 T_{\rm sym}^{\mu\nu}
 =2S\left(p^\mu k^\nu+p^\nu k^\mu-\eta^{\mu\nu}p\!\cdot\!k\right)
 -4R\,ab\,\eta^{\mu\nu},
\end{equation*}
其中
\begin{equation*}
 S\equiv|c_L|^2+|c_R|^2,\qquad
 R\equiv\operatorname{Re}(c_Lc_R^*).
\end{equation*}
有质量矢量场的三个物理极化给出
\begin{equation*}
 \Pi_{\mu\nu}(P)
 =-\eta_{\mu\nu}+\frac{P_\mu P_\nu}{q_V^2}.
\end{equation*}
由二体在壳关系
\begin{equation*}
 p\!\cdot\!k=\frac{q_V^2-a^2-b^2}{2},\qquad
 P\!\cdot\!p=\frac{q_V^2+a^2-b^2}{2},\qquad
 P\!\cdot\!k=\frac{q_V^2-a^2+b^2}{2},
\end{equation*}
先算动量张量的投影：
\begin{align*}
 &\Pi_{\mu\nu}
 \left(p^\mu k^\nu+p^\nu k^\mu-\eta^{\mu\nu}p\!\cdot\!k\right)\\
 &\quad=p\!\cdot\!k
 +\frac{2(P\!\cdot\!p)(P\!\cdot\!k)}{q_V^2}\\
 &\quad=q_V^2-\frac{a^2+b^2}{2}
 -\frac{(a^2-b^2)^2}{2q_V^2}.
\end{align*}
同时
\begin{equation*}
 \Pi_{\mu\nu}\eta^{\mu\nu}=-4+1=-3.
\end{equation*}
因此对初态三个极化平均后的振幅平方为
\begin{align*}
 \overline{|\mathcal M|^2}
 &=\frac13\Pi_{\mu\nu}T_{\rm sym}^{\mu\nu}\\
 &=\frac{q_V^2}{3}
 \left\{
 S\left[2-r_1-r_2-(r_1-r_2)^2\right]
 +12R\sqrt{r_1r_2}
 \right\},
\end{align*}
其中 $r_i=m_i^2c^4/M_V^2=a_i^2/q_V^2$。

精确二体 Lorentz 不变相空间在母粒子静止系为
\begin{align*}
 \int\dd\Phi_2
 &=\frac{1}{8\pi}
 \frac{\sqrt{\lambda(q_V^2,a^2,b^2)}}{q_V^2}\\
 &=\frac{\sqrt{\lambda(1,r_1,r_2)}}{8\pi}.
\end{align*}
将它和上式代回正文\eqnrefs{eq:generic-vector-ff-exact-width-dp92}，并用
$q_V^2=M_V^2/c^2$，得到
\begin{equation*}
 \Gamma_{V\to12}
 =N_c^f\frac{M_V}{48\pi\hbar}\sqrt{\lambda(1,r_1,r_2)}
 \left\{
 S\left[2-r_1-r_2-(r_1-r_2)^2\right]
 +12R\sqrt{r_1r_2}
 \right\},
\end{equation*}
即正文\eqnrefs{eq:generic-vector-ff-exact-closed-dp92}。

正文常用的轻费米子公式定义为有限质量闭式在原点的连续极限。由于
\begin{equation*}
 \lim_{(r_1,r_2)\to(0,0)}\sqrt{\lambda(1,r_1,r_2)}=1,
 \qquad
 \lim_{(r_1,r_2)\to(0,0)}\sqrt{r_1r_2}=0,
\end{equation*}
把 $r_1=r_2=0$ 代入前面的有限质量闭式，其方括号化为 $2S$，从而
\begin{align*}
\lim_{(m_1,m_2)\to(0,0)}\Gamma_{V\to12}
&=N_c^f\frac{M_V}{48\pi\hbar}(2S)\\
&=N_c^f\frac{M_V}{24\pi\hbar}S.
\end{align*}
即
\begin{equation*}
 \boxed{
 \Gamma_{V\to f\bar f}^{(m_f=0)}
 \equiv
 \lim_{(m_1,m_2)\to(0,0)}\Gamma_{V\to12}
 =N_c^f\frac{M_V}{24\pi\hbar}S.}
\end{equation*}
这就是正文\eqnrefs{eq:generic-vector-ff-width-dp91} 所使用的无质量末态子模型；任意有限末态质量均由前面的精确闭式而不是由该极限式描述。

在无质量末态子模型中，对 $W\to\ell\nu$ 代入 $N_c^f=1$、$c_L=g/\sqrt2$、$c_R=0$，得到
\begin{equation*}
 \Gamma(W\to\ell\nu)=\frac{g^2M_W}{48\pi\hbar}.
\end{equation*}
再用 $g=2M_W/v_E$、$v_E^{-2}=\sqrt2G_F^{(E)}$，得到正文\eqnrefs{eq:w-leptonic-width-dp11}。有限 $m_\ell,m_\nu$ 时应回到\eqnrefs{eq:generic-vector-ff-exact-closed-dp92}。同样，在无质量末态子模型中对 $Z\to f\bar f$，
\begin{equation*}
 |c_L^Z|^2+|c_R^Z|^2
 =\frac{g^2}{2c_W^2}\left[(g_V^f)^2+(g_A^f)^2\right],
\end{equation*}
故
\begin{equation*}
 \Gamma(Z\to f\bar f)
 =N_c^f\frac{g^2M_Z}{48\pi c_W^2\hbar}
 \left[(g_V^f)^2+(g_A^f)^2\right].
\end{equation*}
最后用 $M_W=c_WM_Z$ 与\eqnrefs{eq:gf-matching-dp11} 得到正文\eqnrefs{eq:zff-width-tree-dp11}。
\end{derivation}

\begin{derivation}[从\eqnrefs{eq:sm-left-weyl-anomaly-set-dp11}到\eqnrefs{eq:sm-gauge-anomaly-cancellation-dp32}]
先说明为什么规范反常只需要检查正文列出的那些群论迹。对一个左手 Weyl 费米子，三个外部规范玻色子与费米子环相连时，两种相反顶角次序的群因子分别为
\begin{equation*}
 \Tr_R(T^aT^bT^c),
 \qquad
 \Tr_R(T^aT^cT^b).
\end{equation*}
三角图的宇称奇部分在交换后两个外腿时是对称的，因此两种次序相加后只保留
\begin{equation*}
 \Tr_R\!\left[T^a\{T^b,T^c\}\right]
 =\frac12d_R^{abc}.
\end{equation*}
动量积分与正则化决定共同的数值前因子，而一个给定费米子谱是否无规范反常只取决于全部左手 Weyl 场的这些群因子之和是否为零。若其中一个或多个外腿属于 $U(1)_Y$，相应生成元就是超荷乘单位矩阵，于是四类局域反常系数分别化为
\begin{align*}
 [SU(N)]^2U(1)_Y &: \sum_f Y_fT(R_f),\\
 [U(1)_Y]^3 &: \sum_f d(R_f)Y_f^3,\\
 {\rm grav}^2U(1)_Y &: \sum_f d(R_f)Y_f,\\
 [SU(N)]^3 &: \sum_f d_{R_f}^{abc}.
\end{align*}
这里 $d(R_f)$ 表示其余内部指标的多重度。规范理论要求这些群论系数在全部左手 Weyl 场上求和后同时为零。

现在把一代标准模型全部改写成左手 Weyl 场
\begin{equation*}
 Q_L:(\mathbf3,\mathbf2)_{1/6},\quad
 u_R^c:(\bar{\mathbf3},\mathbf1)_{-2/3},\quad
 d_R^c:(\bar{\mathbf3},\mathbf1)_{1/3},\quad
 L_L:(\mathbf1,\mathbf2)_{-1/2},\quad
 e_R^c:(\mathbf1,\mathbf1)_{1}.
\end{equation*}
把共同的三角圈动量积分和规范耦合常数提出以后，定义纯群论反常系数
\begin{align*}
 \mathfrak A_{33Y}
 &\equiv\sum_{\text{左手 Weyl }f}d_2(f)Y_fT_3(R_f),\\
 \mathfrak A_{22Y}
 &\equiv\sum_f d_3(f)Y_fT_2(R_f),\\
 \mathfrak A_{YYY}
 &\equiv\sum_f d_3(f)d_2(f)Y_f^3,\\
 \mathfrak A_{{\rm grav}^2Y}
 &\equiv\sum_f d_3(f)d_2(f)Y_f.
\end{align*}
其中 $d_3,d_2$ 分别是颜色和弱表示维数。利用 $T(\mathbf3)=T(\bar{\mathbf3})=T(\mathbf2)=1/2$，一代标准模型逐项给
\begin{align*}
 \mathfrak A_{33Y}
 &=2\!\left(\frac16\right)\frac12
 +\left(-\frac23\right)\frac12
 +\left(\frac13\right)\frac12=0,\\
 \mathfrak A_{22Y}
 &=3\!\left(\frac16\right)\frac12
 +\left(-\frac12\right)\frac12=0,\\
 \mathfrak A_{YYY}
 &=6\left(\frac16\right)^3
 +3\left(-\frac23\right)^3
 +3\left(\frac13\right)^3
 +2\left(-\frac12\right)^3+1^3=0,\\
 \mathfrak A_{{\rm grav}^2Y}
 &=6\left(\frac16\right)
 +3\left(-\frac23\right)
 +3\left(\frac13\right)
 +2\left(-\frac12\right)+1=0.
\end{align*}
这些等式就是四类反常的群论系数本身为零，而不是只说明它们“与某个零表达式成正比”。
对 $[SU(3)_c]^3$，基本表示与反基本表示的 $d_R^{abc}$ 符号相反，而 $Q_L$ 有两个弱分量，因此
\begin{equation*}
 2d_{\mathbf3}^{abc}+d_{\bar{\mathbf3}}^{abc}+d_{\bar{\mathbf3}}^{abc}
 =(2-1-1)d_{\mathbf3}^{abc}=0.
\end{equation*}
$SU(2)$ 基本表示是伪实表示，局域三次群不变量本身为零。因此一代费米子谱的全部局域规范反常系数均消失，即得到正文\eqnrefs{eq:sm-gauge-anomaly-cancellation-dp32}。正文另行陈述的 Witten 模二条件属于大规范变换的全局拓扑定理，不由这个微扰三角图计算推出。
\end{derivation}






